%%%%%%%% ICML 2026 EXAMPLE LATEX SUBMISSION FILE %%%%%%%%%%%%%%%%%

\documentclass{article}

% Recommended, but optional, packages for figures and better typesetting:
\usepackage{microtype}
\usepackage{graphicx}
\usepackage{subcaption}
\usepackage{booktabs} % for professional tables
\usepackage{multirow}
\usepackage{makecell}
\usepackage{colortbl}
\usepackage{enumitem}
% hyperref makes hyperlinks in the resulting PDF.
% If your build breaks (sometimes temporarily if a hyperlink spans a page)
% please comment out the following usepackage line and replace
% \usepackage{icml2026} with \usepackage[nohyperref]{icml2026} above.
\usepackage{hyperref}


% Attempt to make hyperref and algorithmic work together better:
\newcommand{\theHalgorithm}{\arabic{algorithm}}

% Use the following line for the initial blind version submitted for review:
% \usepackage{icml2026}

% For preprint, use
% \usepackage[preprint]{icml2026}
\usepackage[preprint]{shmicml2026}

% If accepted, instead use the following line for the camera-ready submission:
% \usepackage[accepted]{icml2026}

\usepackage{amsmath}
\usepackage{amssymb}
\usepackage{mathtools}
\usepackage{amsthm}


% if you use cleveref..
% \usepackage[capitalize,noabbrev]{cleveref}

%%%%%%%%%%%%%%%%%%%%%%%%%%%%%%%%
% THEOREMS
%%%%%%%%%%%%%%%%%%%%%%%%%%%%%%%%
\theoremstyle{plain}
\newtheorem{theorem}{Theorem}[section]
\newtheorem{proposition}[theorem]{Proposition}
\newtheorem{lemma}[theorem]{Lemma}
\newtheorem{corollary}[theorem]{Corollary}
\theoremstyle{definition}
\newtheorem{definition}[theorem]{Definition}
\newtheorem{assumption}[theorem]{Assumption}
\theoremstyle{remark}
\newtheorem{remark}[theorem]{Remark}

% Todonotes is useful during development; simply uncomment the next line
%    and comment out the line below the next line to turn off comments
%\usepackage[disable,textsize=tiny]{todonotes}
% \usepackage[textsize=tiny]{todonotes}


%%%%%%%%%%%%%%%%%%%%%%%%%%%%%%%%
% Custom stuff
%%%%%%%%%%%%%%%%%%%%%%%%%%%%%%%%
\usepackage{wrapfig}
\usepackage{amssymb}
\usepackage{enumitem}
\usepackage{etoolbox}
\usepackage[dvipsnames]{xcolor}
\usepackage{graphicx}
\usepackage[capitalize,noabbrev,nameinlink]{cleveref}
\usepackage[breakable]{tcolorbox}
\usepackage{fancyhdr}

\definecolor{lightgray}{gray}{0.90}

\definecolor{cornflowerblue}{rgb}{0.39, 0.58, 0.93}
\definecolor{cornflowerscomplement}{rgb}{0.93, 0.74, 0.39}
\definecolor{cornflowersanalogouspurple}{rgb}{0.47, 0.39, 0.93}
% https://www.canva.com/colors/color-wheel/print/6394ed+7763ed+63d9ed/
\hypersetup{
    colorlinks=true,
    linkcolor=cornflowerblue,
    filecolor=magenta,      
    urlcolor=cornflowersanalogouspurple,%teal,
    citecolor=cornflowerblue,%teal,
    % pdftitle={Accelerated Decoding for Autoregressive LMs via Distillation},
    pdftitle={Multi-Token Prediction via Self-Distillation},
    % pdfpagemode=FullScreen,
    }


% The \icmltitle you define below is probably too long as a header.
% Therefore, a short form for the running title is supplied here:
% \icmltitlerunning{Accelerated Decoding for Autoregressive LMs via Distillation}
\icmltitlerunning{Multi-Token Prediction via Self-Distillation}

\begin{document}

\twocolumn[
  % \icmltitle{Accelerated Decoding for Autoregressive LMs via Distillation}
  \icmltitle{Multi-Token Prediction via Self-Distillation}

  % It is OKAY to include author information, even for blind submissions: the
  % style file will automatically remove it for you unless you've provided
  % the [accepted] option to the icml2026 package.

  % List of affiliations: The first argument should be a (short) identifier you
  % will use later to specify author affiliations Academic affiliations
  % should list Department, University, City, Region, Country Industry
  % affiliations should list Company, City, Region, Country

  % You can specify symbols, otherwise they are numbered in order. Ideally, you
  % should not use this facility. Affiliations will be numbered in order of
  % appearance and this is the preferred way.
  \icmlsetsymbol{equal}{*}

  \begin{icmlauthorlist}
    \icmlauthor{John Kirchenbauer}{umd}
    \icmlauthor{Abhimanyu Hans}{umd}\\
    \icmlauthor{Brian Bartoldson}{llnl}
    \icmlauthor{Micah Goldblum}{cu}
    \icmlauthor{Ashwinee Panda}{umd,tai}
    \icmlauthor{Tom Goldstein}{umd}
  \end{icmlauthorlist}

  \icmlaffiliation{umd}{University of Maryland}
  \icmlaffiliation{cu}{Columbia University}
  \icmlaffiliation{tai}{TogetherAI}
  \icmlaffiliation{llnl}{Lawrence Livermore National Laboratory}

  \icmlcorrespondingauthor{John Kirchenbauer}{jkirchen@umd.edu}

  % You may provide any keywords that you find helpful for describing your
  % paper; these are used to populate the "keywords" metadata in the PDF but
  % will not be shown in the document
  \icmlkeywords{Machine Learning, ICML}

  \vskip 0.3in
]

% this must go after the closing bracket ] following \twocolumn[ ...

% This command actually creates the footnote in the first column listing the
% affiliations and the copyright notice. The command takes one argument, which
% is text to display at the start of the footnote. The \icmlEqualContribution
% command is standard text for equal contribution. Remove it (just {}) if you
% do not need this facility.

% Use ONE of the following lines. DO NOT remove the command.
% If you have no special notice, KEEP empty braces:
\printAffiliationsAndNotice{}  % no special notice (required even if empty)
% Or, if applicable, use the standard equal contribution text:
% \printAffiliationsAndNotice{\icmlEqualContribution}

\begin{abstract}
Existing techniques for accelerating language model inference, such as speculative decoding, require training auxiliary speculator models and building and deploying complex inference pipelines. 
We consider a new approach for converting a pretrained autoregressive language model from a slow single next token prediction model into a fast standalone multi-token prediction model using a simple online distillation objective. 
The final model retains the exact same implementation as the pretrained initial checkpoint and is deployable without the addition of any auxiliary verifier or other specialized inference code. 
Our method produces models that decode more than $3\times$ faster at $<5\%$ drop in accuracy on GSM8K relative to the single token decoding performance of the same checkpoint.
\end{abstract}



\section{Introduction}

Standard language models generate text by predicting tokens one-at-a-time, resulting in very slow throughput when generating many tokens.  This problem is exacerbated in reasoning models, as they often use thousands of tokens during a chain of thought even when ultimately producing a short final response, resulting in a slow (and expensive) user experience.

We consider paradigms for training a language model (LM) to produce multiple tokens in a single forward pass, enabling LMs to generate spans of text with significantly reduced latency and inference cost.  Existing paradigms for multi-token prediction (MTP) rely on variants of the standard cross-entropy loss for next-token prediction (NTP) in which the model output is compared against ground truth texts.  This approach has fundamental limitations that make it difficult or even impossible to generate long, grammatical spans of text.
\begin{figure}[t!]
    \includegraphics[width=\columnwidth]{auto_figures/example_chunk_viz.pdf}
    \caption{Example response to a GSM8K test question from our Qwen3-4B-Instruct-2507 based multi-token prediction model. Decoding is performed using a confidence-adaptive strategy with a threshold of 90\%. Each colored block corresponds to a chunk of tokens produced during a single forward pass and is annotated with its size in tokens which ranges from 1 to 7 in this example. The average chunk size over the entire generation is 3.04.}
    \label{fig:explainer-example-chunks}
\end{figure}

In this paper, we propose an RL-inspired training paradigm in which a student model generates a span of simultaneous token predictions. To avoid the pitfalls of the standard offline objective, the student output is scored by an LM critic/teacher, rather than being scored against a known ground-truth token sequence. By comparing the student's predictions against the next-token suggestions made by the teacher, we produce an on-policy reward signal that enables the student to quickly improve the quality of its multi-token predictions.

After exploring how to implement this training objective in an efficient and scalable way, we turn our attention to inference time sampling strategies.  While our method is capable of producing blocks of $k$ tokens on every forward pass for $k$ as large as 16, 
we can avoid compromising accuracy while still achieving robust speedups
using an adaptive strategy in which tokens are only kept if they have high confidence. 
This strategy naturally enables the model to produce multiple ``easy'' tokens on a forward pass, while focusing its costly single-token passes on ``hard'' tokens that require more computational effort.  This automatic moderation of the speed vs. accuracy tradeoff in various domains like grade school math, instruction following, and open-ended generation allows us to achieve between $2\times$ and $5\times$ acceleration with minimal impact on generation quality.


\subsection{What makes MTP hard?}

A standard NTP language model is trained by minimizing the cross-entropy loss between the softmax probabilities it outputs and the ground truth tokens in a large corpus. 
Under this objective, it learns to model the marginal distributions at each position rather than the joint distribution over full sequences of tokens.
In the NTP context, this does not create an issue at inference time as the autoregressive sampling process allows later token sampling choices to be made in a way that is consistent with previous ones. 
However, in the MTP context, if the distribution for each position in a potential multi-token generation has significant entropy, coherent sampling becomes challenging.  
When multiple tokens are sampled simultaneously from the corresponding marginal distribution, adjacent tokens may be grammatically incompatible.

We can illustrate this problem more clearly using an example. Consider a model trained on the two sentences ``The zookeeper fed the panda bamboo'' and ``The zookeeper fed the lion meat''. If we predict two tokens conditioned on the prefix ``The zookeeper fed the'', then an optimal model should assign equal probability to $\{panda, lion\}$ for the next token, followed by equal probabilities for $\{bamboo, meat\}$.  If we sample both words at the same time, then half the time our zookeeper  feeds the panda meat or the lion bamboo. In the abstract, the LM samples each word independently, and ignores the correlations in the joint distribution between tokens.  For MTP to succeed, we need a training loss, and a sampler, that cares about the {\em joint} distribution.

This problem becomes even more readily apparent in the context of trying to produce much longer generations using a MTP model.
Imagine the extreme case of predicting 100 tokens into the future for a typical sample from an English language training corpus.  Because the text in training sets is (structured) random and unpredictable, the most likely word 100 positions in the future is almost certainly ``the'', as this is the most common word in the English language. Likewise, the most likely word 101 positions into the future is also ``the''. If a model achieves optimal cross-entropy loss at each position, a greedy sampler should predict ``...the the the...'' for far future token positions rather than a useful sentence.  

Our approach overcomes these difficulties by two mechanisms.  First, by using a greedy sampler for the student, we make its output deterministic. We do not need to worry about randomness in the sampler resulting in incompatible adjacent tokens.  Second, because our model output is judged by an LM-based critic using a non-convex loss, degenerate outputs like ``the the the...'' will be rejected.  Under the supervision of a strong teacher, the model must produce valid utterances in the language to achieve low loss, even if the student has a lot of freedom in how to construct those utterances.


\section{Related Work} 

\paragraph{Multi token prediction.}\label{sec:rel-mtp}

An array of recent work has explored architectural and pretraining recipe modifications for building multi-token-prediction language models (MTP LMs). 
For example, \citet{samragh2025your} rely on a pretrained LM's ability to predict multiple tokens in the future, but enhance coherence of the predicted tokens by augmenting the unembedding layer with a new MLP that integrates the previously sampled token's embedding into the current token's representation immediately prior to the unembedding step.
\citet{samragh2025your} also use their MTP tokens as draft tokens in a speculative decoding process, further improving generation quality. 

Based on the observation that a self-distilled training dataset may better align MTP predictions with baseline model predictions, \citet{cai2025fastmtp} use self-distillation with a more traditional cross-entropy loss to finetune a single-layer MTP module. Self-distilled training data has also been used to produce novel MTP modeling approaches like the Parallel Token Prediction models in \citet{draxler2025parallel} and the Jacobi Forcing Models in \citet{hu2025fast}.

We primarily motivate MTP as a means to generate tokens faster than standard sequential decoding without significant loss in quality, but other work instead treats MTP as a component in a larger training objective to improve next token modeling ability \citep{liu2024deepseek}. Finally, other work studies the potential of MTP to address failure modes in long-horizon planning, as it can reduce the degree to which models rely upon the ground-truth prefix provided during teacher-forcing offline training \citep{bachmann2024pitfalls,mahajan2025beyond}. 

\paragraph{Speculative decoding.} While MTP can accelerate token generation, \citet{leviathan2023fast},  \citet{chen2023accelerating}, and \citet{li2025eagle} show that token generation can also be accelerated by using a faster LM to sequentially produce draft tokens that a larger LM can verify in parallel, only accepting tokens that the larger LM would also generate if it had been directly sampled from---a strict but lossless criterion.
This speculative decoding (SD) approach can significantly accelerate token generation by several times depending on the model pairing and text domain. Some recent SD approaches even use a dedicated MTP module as the speculator \citep{cai2024medusa}. However, the actual speedup factors offered by SD in real deployment scenarios have been called into question due to factors not fully considered in academic studies like multi-user batching and the specific datasets considered \citep{liu2025speculative}. 

\paragraph{Language model distillation.} Language modeling capabilities of a teacher model can be transferred to a student through various knowledge distillation objectives~\citep{kim2016sequence}. \citet{agarwal2024policy} perform distillation via feedback from a teacher on sequences generated by the student, ie. ``on-policy distillation'', to avoid distributional mismatch between the student's preferred generations and the static training data. \citet{zhou2023distillspec} apply distillation to align draft models with target models to improve acceptance rates in speculative decoding. 


\paragraph{Online RL and entropy minimization.}\label{sec:rel-ent-min}

Prior work has investigated the role of entropy minimization in both on- and off-policy training scenarios. \citet{wang2021tent} first studied the generalization benefits of direct entropy minimization for image classifiers while modern studies have demonstrated a similar effect in LMs trained with reinforcement learning from verifiable rewards~\cite{agarwal2025unreasonable,shao2024deepseekmathpushinglimitsmathematical}.
\citet{kang2025scalable} exploit the correlation between model confidence and correctness to perform best-of-$n$ sampling at test time. \citet{fu2025deep} leverage the same correlation to re-weight the different roll outs in a pool during online and offline RL, observing consistent performance improvements. 
We build on these results by specifically designing our MTP objective and inference strategies around the correspondence between low entropy predictive distributions and generation quality in reasoning intensive domains.

\section{Methodology}\label{sec:prelim}

In \cref{sec:prelim-ntp} through \cref{sec:prelim-mtpo2} we motivate and define our novel MTP LM training objective in an abstract manner and in \cref{sec:prelim-inference} we discuss how to perform inference using this kind of model. Then in \cref{sec:impl} we discuss specific aspects of our method in more concrete, implementation focused terms with accompanying visualizations.

% \subsection{Notation}\label{sec:prelim-notation}
\subsection{Next Token Prediction}\label{sec:prelim-ntp}
Let $\mathcal{V}$ be the set of $V$ tokens in the vocabulary of our LM, and consider an input sequence $X = (x_1,x_2,...x_N) \in \mathcal{V}^N$. A transformer LM $f_\theta(\cdot)$ parametrized by neural network parameters $\theta$ maps an input $X$ into a sequence of logit vectors $\{\ell_i\in \mathbb{R}^V$\}: 
\begin{equation*}
f_\theta(X) = (\ell_1,\ell_2,..\ell_N) \in \mathbb{R}^{N \times d}.
\end{equation*}
To produce tokens from these logits, we consider a \textit{readout} function $g: \mathbb{R}^{1 \times d} \to \mathcal{V}$ which could be:
\begin{align*}
g(\ell_i) &= \operatorname{argmax}_j(\ell_{ij}) \quad \text{or} \\
g(\ell_i) &= \operatorname{sample}(\operatorname{softmax}(\ell_i)).
\end{align*}
For brevity, in subsequent sections we refer to the $\operatorname{sample}(\operatorname{softmax}(\cdot))$ composition as just ``softmax'' when contrasting it with ``argmax''.

Let $Y$ be the ground truth next tokens at each position in $X$ such that $y_i = x_{i+1}$, and let the predicted next token distribution parametrized by the model be 
\begin{equation*}
P_\theta(x_{i+1}| x_{1:i}) := g(\ell_i) = \operatorname{sample}(\operatorname{softmax}(\ell_i)).
\end{equation*}
If the ground truth next tokens are represented as one-hot vectors denoted $y_i= \{0,1\}^{|\mathcal{V}|}$, then the cross-entropy training objective for a \textit{next token prediction} language model is given as:
\begin{equation}\label{eq:ntp-ce-sum-log}
\mathcal{L}_{NTP} = - \frac{1}{N}\sum_{i=1}^{N} \log P_\theta(y_i | x_{1:i}).
\end{equation}

\subsection{Multi-Token Prediction}\label{sec:prelim-mtp}

In our MTP model, the input $x$ contains prefix tokens, followed by a block of $k-1$ MTP ``mask'' tokens.
On a forward pass, the model predicts $k$ tokens, one at the final prefix position and $k-1$ tokens at each masked position. We can define readouts that operate over multiple logit vectors instead of one:
\begin{align*}
g(l_{i:i+k}) &= \operatorname{argmax}_j(\ell_{i:i+k,j}) \quad\text{or} \\
g(l_{i:i+k}) &= \operatorname{sample}(\operatorname{softmax}(\ell_{i:i+k})).
\end{align*}
Therefore, the \textit{multi-token-prediction} distribution parametrized by the model under the probabilistic readout would be
\begin{equation*}
P_\theta(x_{i:i+k} | x_{1:i}) := g(\ell_{i:i+k}) = \operatorname{softmax}(\ell_{i:i+k}).
\end{equation*}

A naive training loop for MTP can be constructed using a standard offline objective in which all $k$ token predictions are compared to $k$ ground-truth targets, computing the cross-entropy loss. However, in the next section we will introduce our proposal for a MTP training objective that avoids the potential pitfalls with such an approach.

\subsection{``Student Forced'' Online MTP}\label{sec:prelim-mtpo2}

We propose an RL-inspired alternate objective that rewards a coherent joint token distribution in a natural way. 
Consider a strong oracle NTP LM that will serve as our critic/teacher.
Rather than use the ground truth next token distribution as the target distribution, we instead train the student MTP model, $P_{\theta_S}$, to predict $k$ next tokens that a strong \textit{teacher} NTP model, $P_{\theta_T}$, would assign a high likelihood to.  In our above example, an oracle LM would assign a low likelihood (high loss) to ``The zookeeper fed the lion bamboo,'' and so the student trained to minimize the teacher's criticism would not produce this sentence.

As the teacher is a standard NTP model, it does not parametrize the joint token distribution natively. However, by the chain rule of probability, we can use it to score the likelihood of a $k$ token sequence $y' \in \mathcal{V}^k$ proposed by the student. To materialize this sequence, we apply the deterministic readout to the student's logits so that $y':= (y_1,y_2,...y_k) = \operatorname{argmax}(\ell_{i:i+k})$. Then, we compute the likelihood of these tokens under the teacher conditioned on the ground truth prefix according to
\begin{equation}\label{eq:teacher-decomp}
P_{\theta_T}(y' | x_{1:i}) = \prod_{j=1}^{k} P_{\theta_T}(y'_{j} | y'_{1:j-1} \oplus x_{1:i}),
\end{equation}
where $\oplus$ denotes sequence concatenation. 

Finally, we define our new \textit{online} objective as the KL divergence between the likelihood assigned to $y'$ by the student-forced teacher and the student's MTP distribution over its own generation $y'$:
\begin{align}\label{eq:mtpo2}
\mathcal{L}_{MTP} = &- P_{\theta_T}(y' | x_{1:i}) \log P_{\theta_S}(y' | x_{1:i})
% & - H\left( P_{\theta_T}(y' | x_{1:i}) \right). \nonumber
\end{align}
Note that even though it is non-zero in general, we omit the entropy term in the KL involving only $P_{\theta_T}$ as it is constant w.r.t. optimization of the student parameters $\theta_S$.

\paragraph{Advantages of the online approach.}
The objective presented in \cref{eq:mtpo2} has a few desirable properties. First, the student forcing operation means the objective is on-policy. We actually generate a roll out sequence of $k$ tokens from the student model and then compute a ``reward'' based on how likely a strong oracle model thinks they are. Note that, unlike in conventional RL for LMs, our roll outs comprise a single forward pass.

Second, it is also stable across training steps and sequences in the full training corpus. Computing next token likelihoods in service of \cref{eq:teacher-decomp} requires that the teacher model be implemented using the softmax readout function. However, we are not \textit{sampling} from the teacher's distribution, only materializing it. This ensures that the feedback the teacher provides given the same ground truth prefix and student proposal $y' \oplus x_{1:i}$ is deterministic and non-contradictory. 
We expect that an online MTP objective such as ours allows the student model to more efficiently learn to produce a coherent MTP distribution.

\paragraph{``Hard'' teacher distribution.}
Finally, to further enhance the clarity of our supervision signals, we also consider a version of \cref{eq:mtpo2} where the readout function applied to the teacher is also the argmax. This reduces $P_{\theta_T}(y'| x_{1:i})$ to a delta distribution with zero entropy. By construction, computing the loss using a ``hard'' version of the student-forced teacher's distribution requires that the student's entropy drop to zero to minimize it completely. While we do not achieve such full convergence in our experiments, we do observe the desired entropy reduction effect one would expect under this setting and adopt it for our main training runs.

\paragraph{Choice of teacher model.}

While \cref{eq:mtpo2} imposes no particular constraints on the exact model chosen as the teacher or the initialization of the student policy, we believe a natural choice exists. In the special case where $k=1$, during the first step of training, the online MTP objective reduces to a standard NTP knowledge distillation approach; the student need only predict a single token following a given prefix that is likely under the teacher. For this reason, we posit that initializing the student and the teacher using the same pretrained checkpoint, $\theta_T^0 = \theta_S^0$, ensures that the initial loss function is well behaved, especially in the early stage of training. Under this initialization, at step 0 and $k=1$, the student and teacher's predictions are exactly the same and \cref{eq:mtpo2} is 0. Therefore, we will run all experiments under $\theta_T^0 = \theta_S^0$ initialization, but with $k \geq 2$ to allow the student to learn to perform accelerated inference.

\begin{figure*}[t!]
    \includegraphics[width=\textwidth, trim=1cm 4cm 0cm 0cm, clip]{auto_figures/example_tok_masking_viz.pdf}
    \caption{Visual depiction of how a piece of training text is tokenized and masked, including replication of the $k$ ground truth tokens corresponding to each MTP region in which we are actually doing prediction and the corresponding position embedding adjustments required. In this example, the sequence length is 18, the $k$ value is 3, and the number of MTP regions is also 3. Note that while the targets row (\texttt{TgtID}) shown comprises ground truth tokens from the dataset for clarity, under our proposed online training objective, the targets at predicted positions (\texttt{Pred}) are based on the teacher model's feedback, not the ground truth data. \textbf{The masking style and and input replication shown materializes many different MTP problems within a single sequence in parallel, increasing training efficiency.}}
    \label{fig:explainer-tok-masking}
\end{figure*}

\subsection{Inference}\label{sec:prelim-inference}

Having established that the objective function in \cref{eq:mtpo2} may have beneficial properties for MTP LM training, we address how to sample from this type of MTP model at inference time to actually generate sequences of text. 

The training objective \cref{eq:mtpo2} uses the deterministic argmax readout function to materialize the student's $k$-token rollout $y'$. However, at inference time, we must also choose a method for sampling from the student. We note that in our simple MTP formulation, $P_{\theta_S}(y' | x_{1:i})$ does not explicitly parametrize the product distribution over $\mathcal{Y}^k$ and instead produces $P_{\theta_S}(y'_j | x_{1:i}), \quad j \in [1,..k]$. Therefore, under the softmax readout, the $k$ tokens would actually be sampled individually. If the distribution has any entropy---meaning the student is not perfectly certain about the individual tokens at each position---then the sampled tokens may not be mutually compatible.

In principle, this issue is resolved simply by using the greedy argmax readout function at test time. Because the model is trained using a deterministic argmax, a well trained MTP model should yield argmax tokens at each position that are compatible with their argmax neighbors. 
In principle, the only thing that matters for such a model is which token is assigned the largest logit value rather than the entropy of the model's entire output distribution.
We will determine empirically whether or not argmax tokens produced by a student model trained using~\cref{eq:mtpo2} are actually coherent in practice. Later, we will see that despite the fact that it is not strictly necessary, high generation coherence (low perplexity) appears to be closely correlated with high token confidence (low entropy) in our experiments, and we will exploit this fact to create a simple adaptive sampler.

\section{Implementation Details}\label{sec:impl}

In \cref{sec:prelim}, we introduced our method in an abstract manner. Now, in \cref{sec:impl-init} through \cref{sec:impl-confadapt} we provide details about the actual implementation of our training and inference scheme. \Cref{fig:explainer-tok-masking,fig:explainer-block-mask-roll,fig:explainer-block-mask-vark} are rendered in support of components of our implementation that are most succinctly described using visuals such as attention patterns and input tokenization and masking.

\subsection{Initialization}\label{sec:impl-init}

To minimize training requirements, we initialize the student MTP policy and teacher NTP model from the same set of strong pretrained weights from a NTP model. During training, we keep the teacher's parameters frozen while allowing all parameters in the MTP model to freely train.  We randomly initialize a special \verb|<MTP>| token (as shown in \cref{fig:explainer-tok-masking}) in the embedding and un-embedding matrices of the MTP LM according to the mean and variance of the existing pretrained embeddings.

\subsection{Tokenization and Masking}\label{sec:impl-tok-mask}

The training objective proposed in \cref{eq:mtpo2} is computed over a specific position $i$ in a ground truth sequence $X$ and the next $k$ tokens that the student predicts to follow the prefix $x_{1:i}$. For any sequence $X$ of length $N$, there are $N-1$ unique prefixes that can be materialized as input for our online objective. To train efficiently, we need an attention masking scheme that allows us to supervise multiple prediction spans within a single sequence in parallel.

In \cref{fig:explainer-tok-masking}, we present a visual depiction of how input text is tokenized into a sequence, masked, and aligned for our model so that multiple supervised spans can be computed in parallel. To compute the $k$ token prediction $y'$ of the student at a given position in the tokenized input sequence, $k-1$ ``MTP masks'' are inserted directly following the last token of the prefix that will be provided as input to the model for this specific MTP problem identified by position $i$.

Another important detail shown in \cref{fig:explainer-tok-masking} is the appearance of ground truth tokens following the first MTP region. Suppose we make MTP predictions at position $i$, followed by another prediction at position $i'>i.$ This later prediction must have access to all ground truth tokens in the span $x_{0:i'}$, which includes the tokens $x_{i:i'}$ between $i$ and $i'$.  We place the ground truth tokens $x_{i:i'}$ immediately after the first MTP region. We also modify the position embeddings (e.g. RoPE features) so these tokens have their original position features (ignoring the presence of the MTP mask tokens) applied to them, though they now appear later in the context window.

\subsection{Blocked Attention}\label{sec:impl-block}
\label{sec:block}

We visualize examples of our attention masks in \cref{fig:explainer-block-mask-roll,fig:explainer-block-mask-vark}.
We design the mask such that for all positions corresponding to ground truth tokens outside of the MTP regions, the mask has the standard causal structure; each ground truth token attends to all upstream prefix tokens, skipping over MTP spans. Each MTP token attends to upstream ground truth tokens, in addition to upstream tokens within its own MTP span (but not other spans). 

We note that in the masked regions where MTP is performed, the use of causal attention, rather than bidirectional or full self-attention, is not strictly necessary---whereas it \textit{is} for standard NTP training using offline cross-entropy (see \cref{app:ablations} for a discussion). However, since using a bidirectional mask in the MTP regions would represent an additional distribution shift for the pretrained student model, would also preclude the use of standard causal attention at inference time, and only provides marginal lift in our ablations (\cref{tab:abl-l3-suprv-evals-pt1}) we choose to train our models with blocked but otherwise standard causal attention.

\begin{figure}[h!]
    \includegraphics[width=0.49\columnwidth]{auto_figures/MTP_block_mask_with_Seq_Len=18_Prefix=4_k=3_Offset=0.png}
    \includegraphics[width=0.49\columnwidth]{auto_figures/MTP_block_mask_with_Seq_Len=18_Prefix=4_k=3_Offset=-2.png}
    \caption{Visualization of an attention mask with rolling offsets; offset 0 on the left and -2 on the right for a fixed $k=3$. \textbf{Randomized offsets enable supervision on problems with many different prefix lengths during the same training run.}}
    \label{fig:explainer-block-mask-roll}
\end{figure}
\begin{figure}[h!]
    \includegraphics[width=0.49\columnwidth]{auto_figures/MTP_block_mask_with_Seq_Len=12_Prefix=4_k=3_Offset=0.png}
    \includegraphics[width=0.49\columnwidth]{auto_figures/MTP_block_mask_with_Seq_Len=12_Prefix=2_k=5_Offset=0.png}
    \caption{Visualization of an attention mask with variable $k$ masking showing $k=3$ and $k=5$ with a fixed offset of 0. \textbf{Randomized $k$ values enable supervision on MTP problems of many different sizes during the same training run.}}
    \label{fig:explainer-block-mask-vark}
\end{figure}

\subsection{Random Offsets and $k$ Values}\label{sec:impl-randomize}

We randomize the location of the MTP spans so that the model can learn to make MTP predictions starting from any position.
\cref{fig:explainer-tok-masking,fig:explainer-block-mask-roll} illustrate how we achieve this using an offset when masking tokens in the input and blocking the attention mask. We also experiment with randomizing the MTP span length $k$. As shown in \cref{fig:explainer-block-mask-vark}, we parametrize our attention mask generator to allow dynamic adjustment of position, length, and number of MTP spans.

\subsection{KV Cache Management}\label{sec:impl-kvcache}
 
During inference, our MTP LM appends mask tokens to the prefix before the forward pass. These tokens should be removed from the KV cache after they are used for prediction.  In our implementation, the first prediction pass on a sequence generates $k$ KV pairs, one for each token predicted.  After this pass is complete, the mask tokens are removed from the prefix and the corresponding KV values are popped off the top of the cache. Then, $2k-1$ tokens are appended to the prefix; the $k$ newly generated tokens and $k-1$ new mask tokens for the next forward pass.  We proceed in this way, popping stale values off the cache and adding the new tokens and masks before every forward pass. This scheme allows us to save KV cache space while also using a standard causal mask at inference time (rather than a blocked mask like in~\cref{fig:explainer-block-mask-roll}) which is typically more efficient in most transformer implementations.

\subsection{Static vs. Adaptive Decoding Strategies}\label{sec:impl-confadapt}
Our informal analysis of the proposed objective \cref{eq:mtpo2} in the preliminaries section suggests that there may be a correlation between the confidence of the student model on its predictions and the quality of the sequence emitted.
In \cref{fig:corr-analysis} we present empirical evidence of this relationship from an early experiment. In response to this confirmation we devised a simple \textit{confidence-adaptive} scheme (\textit{ConfAdapt}) that allows us to set the $k$ value dynamically at each decoding step based on this confidence heuristic. Given a threshold value $\tau$, such as $90\%$, the CA strategy finds the greatest index $k' \in (1,k_{max})$ such that $\operatorname{argmax}_j(\ell_{ij}) > \tau, \forall i \in (1,k')$, and uses that index as the $k$ value for this decoding step. This creates a dynamic range of $k$ values during a single generation which we compute statistics over to report average ``acceleration factors''.

While we experiment with randomizing the value of $k$ at each step during the training process, we do not directly train for any sort of dynamic decoding procedure where $k$ varies at each generation step during a long generation. However, to our surprise, the MTP models we train turn out to be amenable to the simple dynamic decoding strategy described above, and it achieves pareto-optimal results when compared to static strategies.
\begin{figure*}[t!]
    \includegraphics[width=0.441\textwidth, trim=0cm 0cm 2.1cm 0cm, clip]{auto_figures/dynamics_daint_prod_ift_mask_fix_1N4n_9d30cad5_gsm8k_cot_retrorec_scatter_only.pdf}
    \includegraphics[width=0.548\textwidth, trim=0cm 0cm 0cm 0cm, clip]{auto_figures/dynamics_daint_prod_ift_q3-4b_1N4n_16cdce0f_gsm8k_cot_retrorec_scatter_only.pdf}
    \caption{The performance of our (\textbf{Left}) L3.1-8B-Magpie based MTP LM and (\textbf{Right}) Qwen3-4B-Inst-2507 MTP LM evaluated on the GSM8K benchmark after $\sim$100k steps of training. Performance tradeoff is visualized by plotting the effective $k$ value or ``Acceleration Factor'' versus the Accuracy on the benchmark. More detailed plots showing accuracy and acceleration as a function of training steps for both models are provided in \cref{fig:dynamics-flagship-l3-gsm,fig:dynamics-flagship-q3-gsm}. \textbf{We observe that the adaptive decoding strategies achieve pareto-optimal tradeoffs between generation speed and response quality for both models.}}
    \label{fig:dynamics-flagship-l3-q3-gsm-scatter}
\end{figure*}

\section{Experiments}\label{sec:exps}

We adapt pretrained NTP LMs into MTP LMs and evaluate their performance on both mathematical reasoning benchmarks and open ended knowledge intensive generative tasks.


\subsection{Setup}\label{sec:exps-setup}

We describe all aspects of our training and evaluation setup in detail in \cref{app:impl-exp-details} but we succinctly note the most critical design parameters here to support the main experimental results. The pair of initial checkpoints used are Llama-3.1-8B-MagpieAlign-SFT-v0.1 (``L3.1-8B-Magpie'' hereafter, \citet{xu2024magpie}) and Qwen3-4B-Instruct-2507 (``Qwen3-4B-Inst-2507'' hereafter, \citet{yang2025qwen3}). In order to target a benchmark involving reasoning traces where inference efficiency is particularly useful in practice, we tune our MTP LMs on a dataset of synthetic grade school math (GSM) problems called MetaMathQA \citep{yu2023metamath}.

Based on the results of preliminary experiments, unless stated otherwise, we train with randomized block mask offsets and randomized $k$ values in the range $[2,16]$; this implicitly sets a reasonable choice of $k_{max}$ at 16 during inference using the \textit{ConftAdapt} scheme. We also use the ``hard teacher'' variation on \cref{eq:mtpo2} where the readout function for the teacher NTP model is the argmax rather than softmax as early experiments suggested that this lowered the student model's output entropy more rapidly. Under $k$ value randomization, in expectation the $k$ value is closer to half of the max. Therefore, over the course of the $\sim 100$k steps of training our models see a total of $\sim500$M supervised tokens (see \cref{app:impl-exp-details} for this calculation).

To assess the performance of our main models we select the following set of \textit{generative} benchmarks to test the technical and long form generation capabilities of our MTP LMs: GSM8K, AIME25, GPQA, BBH, IFEVAL, and CNN DailyMail. We describe the exact evaluation parameters for all datasets in \cref{app:impl-exp-details}. 
However, since we find the results on GSM8K to be both the strongest and also indicative of all major trends, we focus on that task in the main body of this paper.

\textbf{GSM8K CoT Fewshot} is evaluated with an additional ``Lets think step by step'' suffix to the prompt, with output truncation after emission of \verb|``Q:''| or any of the model's stop tokens. We report the ``flexible-extract'' accuracy metric and include its Std. Error in tables.

\subsection{Performance vs. Acceleration}\label{sec:exps-perf-models}

In \cref{fig:dynamics-flagship-l3-q3-gsm-scatter} we summarize the performance of our L3.1-8B-Magpie based and Qwen3-4B-Inst-2507 based MTP models after $\sim 100,000$ steps of training on MetaMathQA. Results under the \textit{ConfAdapt} scheme are visualized on the same axes as the \textit{Static} scheme by computing the average $k$ value that was chosen by the heuristic across all decoding steps over all the generations in the evaluation and include its Std. Error in tables. 

Under the \textit{Static} scheme, as $k$ increases the acceleration factor grows and the accuracy decreases. Similarly, under the \textit{ConfAdapt} scheme, as the confidence threshold $\tau$ drops, the acceleration factor increases while accuracy decays, albeit more slowly. When compared against \textit{Static} $k=1$ decoding, using the \textit{ConfAdapt} scheme at threshold of $90\%$ yields an acceleration factor of more than $3\times$ while decreasing accuracy for the L3.1-8B-Magpie by less than $3\%$; the Qwen3-4B-Inst-2507 achieves $3\times$ speeds at a slightly higher cost of a $7\%$ drop in accuracy. At more aggressive settings using more permissive confidence thresholds, the acceleration factors are as high as $5\times$ but cause a more significant reduction in accuracy.

\subsection{Performance Across Tasks}\label{sec:exps-perf-tasks} 

In \cref{tab:math-evals,tab:general-evals}, we present the same pair of models but evaluated over a larger set of 6 evaluation tasks. 
We observe that even though both models were trained on only MetaMathQA, transfer learning occurs.
As the two initial models have different strengths before MTP adaptation, both at step 0 and after training is complete, their performance under \textit{Static} $k=1$ decoding varies considerably. However, we generally observe that larger average acceleration factors occur on the GSM8K and BBH benchmarks, and that for the more open-ended generative tasks like CNN DailyMail, lower acceleration factors are yielded for all adaptive decoding configurations.


\subsection{Ablations}\label{sec:exps-abls}

In addition to the two main models that we showcase above, in the Appendix we include results for ablations where we change the training dataset and where we modify the various parameters of our MTP training objective. These details and results are presented in \cref{app:ablations} and inform our choice of settings for our primary experiments. Specifically in terms of supervision style, we can summarize the results shown in \cref{tab:abl-l3-suprv-evals-pt1,tab:abl-l3-suprv-evals-pt2} as supportive of our choice of hyperparameters and objective configuration for the main models. We observe the best overall performance when we utilize the hard teacher distribution under the argmax readout as targets, when we randomize the $k$ values during training, when we use causal masking across the MTP regions, and when there is no auxiliary loss term for NTP on prefix tokens.

\subsection{Limitations}\label{sec:limitations}

While we present a series of ablations to demonstrate the efficacy of our proposed objective, early experiments and the outcomes of certain ablations suggest that there are many avenues for improvement. For example, while we observe non-trivial lift when using randomized $k$-values rather than static ones during training, there are other ways we could have performed the comparison that might be more compelling. We also explored utilizing a non-static curriculum over $k$ values but those experiments were inconclusive. We also considered directly penalizing the entropy of the student model to collapse its distribution faster. However, we expect that other approaches such as a second stage of more traditional online RL with complete multi-step roll outs from our trained MTP LMs might prove more effective than the penalty term approach. Finally, the pair of datasets and limited training budgets we utilize yield surprisingly strong results but we expect that scaling up the recipe (dataset, model size, training compute) could all help produce more performant versions of our models.
We release our training code and model checkpoints so that future work can explore some of these underexplored directions.

\section{Conclusion}\label{sec:conc}

Our methodology produces MTP LMs with accelerated decoding abilities using a straightforward recipe and modest computational resources.
While techniques in the speculative decoding literature achieve similar levels of acceleration under certain conditions, they require the training of additional speculator modules and the development and maintenance of complex inference implementations to handle the proposal and verification steps. Our approach is \textit{much} simpler but achieves similar results.
More broadly, our results suggest that decoding acceleration need not rely solely on inference-time mechanisms. By absorbing some complexity into training the model, our MTP strategy provides a complementary and largely orthogonal axis to existing accelerated decoding methods.

\paragraph{Directions for Future Research.}
We conclude by proposing a number of avenues for future work that we believe can be explored without significant computational expenditure. 
  
% enumitem interface
% \begin{enumerate}[label = (\roman*)]
\textbf{(i)} Reducing the entropy of our MTP LM to zero is the goal of the training process, but our training process may not be the best procedure for reducing the entropy. Future work could explore other strategies for decreasing output entropy. We expect that as entropy in the output distribution approaches zero, we can conceivably produce many more tokens during each forward pass while maintaining coherence.

% \item Our goal is to model the distribution of text, which is actually high entropy, but we would like to % my guess at a compl for this stub...
\textbf{(ii)} Even under accelerated decoding, our goal is to still faithfully model the highly entropic distribution of natural text, which is in tension with our acceleration scheme's implicit low entropy requirements. Some simple decoding strategies that would force non-determinism purely at test time are possible (e.g. sampling from the softmax whenever $k=1$), but future work could incorporate a source of randomness more naively into the model to increase roll out diversity.

\textbf{(iii)} Our models could be used for self-speculative decoding. In this work we specifically avoid incorporating a verification pass to explore whether or not fast generation of $k$ tokens in a single shot is even feasible. However, in principle, the proposals made by our MTP model could be used for self-speculative decoding. The MTP model could verify its own predictions during each subsequent forward pass to achieve more modest but lossless speedups. We expect that our KV caching strategy is amenable to an additional step where previously emitted tokens are verified before committing them to the cache (see \cref{sec:impl-kvcache}).

\textbf{(iv)} Our simple MTP implementation involves using a trainable mask token to extend the sequence passed into the model by $k$ positions, creating additional slots for computation within the model, but other MTP architectures are possible (discussed in \cref{sec:rel-mtp}). In principle, adding small auxiliary prediction heads or modules can enable MTP and it is possible that these less computationally intensive implementations are as readily trainable and performant as our proposal, even under the same objective.

\textbf{(v)} Our online distillation objective can be viewed as a simple type of RL, but a full-fledged policy optimization approach may yield better results. For simplicity and parallelizability, during training we use a single forward pass to produce just one block of $k$ tokens at each masked region within a training sample. However, a more long horizon approach is possible where the student MTP model rolls out its own full sequence comprising many chunks of $k$ tokens over multiple forward passes that is scored using a more traditional reward signal. This could result in a beneficial collapse in entropy that improves output quality and acceleration.
% (see \cref{sec:rel-ent-min})
% \end{enumerate}


% Acknowledgements should only appear in the accepted version.
\section*{Acknowledgements}

This work was supported by DARPA TIAMAT, the NSF TRAILS Institute (2229885), and Coefficient Giving. Computing resources for this project were supported in part by the Swiss National Supercomputing Center (CSCS).  Prepared in collaboration with Lawrence Livermore National Lab (LLNL) under Contract DE-AC52-07NA27344 and supported by the LLNL-LDRD Program under Project No. 24-ERD-010 (LLNL-CONF-2015543).

We would also like to acknowledge Jonas Geiping and Sean McLeish for helpful discussions throughout the research process on infrastructure and tooling for experimentation on the public supercomputers at CSCS and LLNL.

\bibliography{references}
\bibliographystyle{icml2026}

%%%%%%%%%%%%%%%%%%%%%%%%%%%%%%%%%%%%%%%%%%%%%%%%%%%%%%%%%%%%%%%%%%%%%%%%%%%%%%%
%%%%%%%%%%%%%%%%%%%%%%%%%%%%%%%%%%%%%%%%%%%%%%%%%%%%%%%%%%%%%%%%%%%%%%%%%%%%%%%
% APPENDIX
%%%%%%%%%%%%%%%%%%%%%%%%%%%%%%%%%%%%%%%%%%%%%%%%%%%%%%%%%%%%%%%%%%%%%%%%%%%%%%%
%%%%%%%%%%%%%%%%%%%%%%%%%%%%%%%%%%%%%%%%%%%%%%%%%%%%%%%%%%%%%%%%%%%%%%%%%%%%%%%
\newpage
\appendix
\onecolumn

\section{Appendix}

\subsection{Extended Implementation and Experimental Details}\label{app:impl-exp-details}

\paragraph{Stopping criteria.}

Similar to KVCaching, the natural generalization of stop-token handling, eg. EOS tokens, from the NTP setting to our MTP model is straightforward but still requires explanation. 

During training, since we primarily consider instruction following datasets with variable lengths, we have to choose how to train the model to emit an EOS token. Our convention is to treat the first MTP region in which an EOS appears in the ground truth sequence as the determiner of the last MTP region for this sample in the training data. Then, we pad out the sequence with more EOS tokens to fill this last $k$ token region. Beyond this region, the rest of the context window is filled with padding and ignored completely by the computation and objective. 

At inference time, we therefore hope that the model has then learned to emit at least one EOS token in the appropriate point in the sequence. During generation, if any of the tokens emitted by the MTP model in a given $k$ token decoding step are stop-tokens, we simply call this the final step of generation and terminate, trimming off any tokens predicted after the first EOS. We find that this simple setup provides the MTP model the ability to reliably generate long sequences that truncate at logical positions that are also not exact multiples of $k$.

\paragraph{Pretrained models.}

Initially we prototyped our approach using the Llama-3.2-1B base checkpoint \citep{grattafiori2024llama} to enable low cost, rapid iteration of our design. Then we upgraded to the Llama-3.1-8B base checkpoint \citep{grattafiori2024llama} to improve overall performance while exploring loss variations and training hyperparameters (\cref{fig:corr-analysis} is derived from experiments with that model). Finally, based on intuitions that post-trained models generally have lower entropy in their NTP distribution due to training on narrow domains using objectives and datasets that empirically reduce a model's entropy, we switched to post-trained initial checkpoints for the final set of experiments and ablations. The pair of post-trained initial checkpoints used are Llama-3.1-8B-MagpieAlign-SFT-v0.1 (``L3.1-8B-Magpie'' throughout, \citet{xu2024magpie}) and Qwen3-4B-Instruct-2507 (``Qwen3-4B-Inst-2507'' throughout, \citet{yang2025qwen3}).

\paragraph{Finetuning data.}

In order to target benchmarks involving reasoning traces where inference efficiency is most useful in practice, we train our MTP LMs on a dataset of synthetic grade school math (GSM) problems called MetaMathQA \citep{yu2023metamath}. The data is a synthetically generated series of GSM problems with reasoning traces and answers built by taking an existing set of previously released, hand-written questions, and augmenting them using a large API based model. We note that this dataset is based on the train split of the official GSM8K dataset only, the official test set questions used for benchmarking are not included or used as seed questions. For our experimentation, we also split the 395k raw rows into subsets for training and validation that are stratified by the unique set of seed questions such that augmentations of a single seed question fall into the train set, or the validation set, but not both.

In an ablation, we also consider a more general set of finetuning data: Magpie-Pro-300k-v0.1 and Magpie-Reasoning-150k (``Magpie'' hereafter, \citet{xu2024magpie}). These are also synthetically generated datasets covering a wide range of instruction following and reasoning tasks, but we specifically choose it due to the release of a publicly released model checkpoint that was post-trained directly on these datasets (the aforementioned L3.1-8B-Magpie). We hypothesize that the affinity between the pretrained initialization of the student policy and its ability to rapidly learn the MTP behavior may be linked and so we compare the impact of this training data choice, holding all other settings constant.

\paragraph{Chat Templating.}

Both of the main pretrained models we use come equipped with a chat template. However, preliminary experiments indicated that it was difficult to ensure that our training data and evaluation data was prepared in such a way that the model responded accurately under standard $k=1$ decoding at the end of training. As a result, we experimented with various combinations before landing on the following configuration. For the L3.1-8B-Magpie model, we apply the chat template included with its tokenizer for all training and evaluation datasets as this performed well. However, for Qwen3-4B-Inst-2507, after experimenting with its chat template and observing poor results, we settled on preparing the training and evaluation data by joining ``input'' and ``response'' pairs using just a simple ``\verb|\n\n|'' string and prepending a BOS token to every full input. In all tables, the setting for the model is noted under its name.

\paragraph{Hyperparameters.}

In addition to the initialization and training datasets, we adopt a standard set of hyperparameters for elements of the training pipeline not specific to our MTP implementation and objective. We train our models using the AdamW optimizer using PyTorch default settings and a learning rate schedule of 2000 steps of warm-up followed by a constant peak learning rate of $1e-5$; our goal is to keep the model in a ``finetuning'' regime where the student's initial capabilities and behaviors are maintained whilst the MTP abilities are acquired. All parameters in the student model are made trainable, while the teacher model remains frozen.

As our primary training dataset, MetaMathQA, contains relatively short samples (avg. $\sim 220$ tokens under the Llama 3 tokenizer), we cap the overall sequence length (training context window, $N$) at 160 tokens for all MetaMathQA runs. For ablations using the Magpie datasets whose average lengths are longer, we use a sequence length of 1024 tokens. For experiments with MetaMathQA, we use a global batch size of 128 sequences across all devices training in data parallel, and for the Magpie dataset experiments we use a global batch size of 16. While not perfectly equivalent, the number of tokens per optimization step is therefore within the same order of magnitude for both experiments. Based on a desire to balance the supervision density of our objective with the length of prefixes that the model is exposed to, we set the number of MTP regions $M$ in a given sequence during training to be $M=N/(2k)$ as illustrated in \cref{fig:explainer-tok-masking} (we use the largest $k$ value 16 as the value when choosing $M$, eg. $M=N/(2k_{max})$).

For implementation simplicity, while the offset and $k$ value are randomized during training $M$ is left fixed meaning that all sampled values of offset and $k$ produce a lower number of supervised positions per (eg. the sum of the ``Pred:'' row in \cref{fig:explainer-tok-masking}) than the maximum which is when the offset is 0 and the sampled $k$ is the maximum. We use a batch size of 128 sequences for the models trained on MetaMathQA, and a batch size of 16 for the models trained on the Magpie datasets. While not scaled exactly, due to the $> 5\times$ larger sequence length for the latter, the reduced batch size still achieves in the same order of magnitude number of tokens per optimization step. Over the course of approximately 100k training iterations, for the MetaMathQA models based on the setting of $M=N/(2k_{max})$, this equates to an upper bound of 1B supervised tokens if assuming a fixed value $k=k_{max}$ and an offset of 0. Since $160 /(2*16) = 5$ regions, we have $100e3 * 128 * (5 * 16) = 1.024$B tokens. For the actual runs under randomization, in expectation the $k$ value is closer to half of the max, so our models only see a total of 500M supervised tokens. 

As illustrated in \cref{fig:explainer-tok-masking}, more tokens of raw input are required as input at each step than are actually supervised if $k$ is randomized following the static $M$ setting described above. As there are about 50M tokens available in the MetaMathQA dataset assuming the Llama 3 tokenizer and a truncation length of 160, this results in $\sim35$ epochs of training for the L3.1-8B-Magpie based model. While early experiments indicate that overfitting is possible at more extreme durations, due to the offset and $k$ value randomization strategies, the number of unique batches of prefix and MTP prediction problems in a dataset with $D$ rows is actually $>D$. These additional views of the same sample---approximately proportional to the unique combinations of block offset and $k$ possible under a given set of hyperparameters---likely make the effective number of epochs on the same exact batch of prepared  and masked inputs much lower, if not close to just 1. We believe that the result of transfer learning from MetaMathQA to benchmarks beyond GSM8K also suggest that overfitting isn't necessarily an issue in our specific experimental settings.

The main models containing 8B and 4B parameters are both trainable at the settings described using one node of 4$\times$GH200 GPUs using FSDP based model and data parallelism. The training runs to 100k steps take approximately 24-36 hours of continuous runtime depending on the model. Each benchmark evaluation can be performed on a single GH200 GPU and runtime varies depending on the specific settings of $k$, the size and average response length for the benchmark itself.


\paragraph{Evaluation metrics and benchmarks.}

During early experimentation, we used our held out validation subset of MetaMathQA to track metrics like confidence and quality of the student roll outs under the teacher (as in \cref{fig:corr-analysis}). Then, for a comprehensive evaluation of our main models we selected a series of mathematical and general instruction following benchmark datasets from the Eleuther LM Evaluation Harness~\citep{eval-harness}. We specifically chose the following set of \textit{generative} benchmarks to test the technical and long form generation capabilities of our MTP LMs: GSM8K, AIME25, GPQA, BBH, IFEVAL, and CNN DailyMail. We refer the reader to the harness repo and each respective dataset's documentation for a detailed descriptions of each but we list the specific used settings in the evaluation harness below.

\textbf{GSM8K CoT Fewshot} is evaluated with an additional ``Lets think step by step'' suffix to the prompt, with output truncation after emission of \verb|``Q:''| or any of the model's stop tokens. We report the ``flexible-extract'' accuracy metric and its Std. Error.

\textbf{AIME25} is evaluated with output truncation after emission of any of the model's stop tokens. We report the ``exact-match'' accuracy metric and its Std. Error.

\textbf{GPQA Main CoT Fewshot} is evaluated with output truncation after emission of any of the model's stop tokens. We report the ``flexible-extract'' accuracy metric and its Std. Error.

\textbf{Big Bench Hard CoT Fewshot} is evaluated with output truncation after emission of \verb|``\\n\\n''|, \verb|``Q:''|, or any of the model's stop tokens. We report the exact-match accuracy metric and its Std. Error under the ``get-answer'' extraction method.

\textbf{IFEval} is evaluated with output truncation after emission of \verb|``Q:''| or any of the model's stop tokens. We report the ``Prompt-level'' exact-match accuracy metric and its Std. Error.

\textbf{CNN DailyMail Summarization} is evaluated with output truncation after emission of any of the model's stop tokens. We report the ROUGE-L performance metric and its Std. Error.


\subsection{Extended Results}

\subsection{Quality vs. Confidence Correlation}\label{sec:conf-qual-corr}

\begin{figure}[h!]
    \centering
    \includegraphics[width=0.6\textwidth, trim=0cm 0cm 0cm 0cm, clip]{auto_figures/corr_SS_Stud_Gen_top1_confs_vs_SS_Stud_Gen_Loss_largeN_k3_ss_steps_117000-165500.png}
    \caption{The correlation between student confidence and generation quality for a preliminary experiment where we train a Llama-3.1-8B Base model on generic webtext under our MTP objective. We evaluate the quality of its completions to prefixes in a validation set (loss under the teacher model, lower is better) and correlate this against the student's confidence in its own generations (avg. probability assigned to top token in all $k$ positions). \textbf{The correlation between confidence and quality allows us to design a simple adaptive decoding strategy that limits the number of tokens sampled if the model is not sufficiently confident at any point during the generation process.}}
    \label{fig:corr-analysis}
\end{figure}

\subsection{Performance over the course of training.}
In \cref{fig:dynamics-flagship-l3-gsm,fig:dynamics-flagship-q3-gsm} we present the evaluation performance and acceleration on the GSM8K benchmark over the course of $\sim 100,000$ training iterations for our L3.1-8B-Magpie based and Qwen3-4B-Inst-2507 based MTP models respectively. We find that when evaluating the models under the \textit{Static} scheme at various values of $k$, performance is inversely proportional to the value of $k$. We also observe that both performance and acceleration stabilize by the halfway point.
Note that we report an initial acceleration factor $< k$ for Static schemes in some cases because when untrained, the model often emits a stop token within the first $k$ token window, truncating the sequence.

The \textit{ConfAdapt} scheme is visualized on the same axes as the Static scheme by computing the average $k$ value that was chosen by the heuristic across all decoding steps over all the generations in the evaluation. 
% We present the combinations of two settings each for the confidence threshold and $k_{max}$. 
We present the results across a sweep of confidence thresholds for each model and benchmark. 
We observe that the ConfAdapt schemes are pareto optimal, achieving significant acceleration factors of $>3\times$ without degrading accuracy very much. At more aggressive settings using more permissive confidence thresholds, the acceleration factors are as high as $5\times$ but cause a more significant reduction in accuracy.
\begin{figure*}[h!]
    \includegraphics[width=\textwidth]{auto_figures/dynamics_daint_prod_ift_mask_fix_1N4n_9d30cad5_gsm8k_cot_retrorec_static_only.pdf}
    \includegraphics[width=\textwidth]{auto_figures/dynamics_daint_prod_ift_mask_fix_1N4n_9d30cad5_gsm8k_cot_retrorec_ca_only.pdf}
    \caption{(Left) Accuracy and (middle) acceleration dynamics for L3.1-8B-Magpie evaluated on GSM8K over the course of training. (Right) The pareto frontier of accuracy versus acceleration based on the final model checkpoint evaluated across all decoding strategies. We find that performance and acceleration stabilize around 50k training steps, and the adaptive decoding strategies achieve pareto-optimal tradeoffs between generation speed and response quality.}
    \label{fig:dynamics-flagship-l3-gsm}
\end{figure*}
\begin{figure*}[h!]
    \includegraphics[width=\textwidth]{auto_figures/dynamics_daint_prod_ift_q3-4b_1N4n_16cdce0f_gsm8k_cot_retrorec_static_only.pdf}
    \includegraphics[width=\textwidth]{auto_figures/dynamics_daint_prod_ift_q3-4b_1N4n_16cdce0f_gsm8k_cot_retrorec_ca_only.pdf}
    \caption{(Left) Accuracy and (middle) acceleration dynamics for Qwen3-4B-Inst-2507 evaluated on GSM8K over the course of training. (Right) The pareto frontier of accuracy versus acceleration based on the final model checkpoint evaluated across all decoding strategies. We find that performance and acceleration stabilize around 50k training steps, and the adaptive decoding strategies achieve pareto-optimal tradeoffs between generation speed and response quality.}
    \label{fig:dynamics-flagship-q3-gsm}
\end{figure*}

\begin{figure*}[t!]
    \includegraphics[width=0.435\textwidth, trim=0cm 0cm 2.1cm 0cm, clip]{auto_figures/dynamics_daint_prod_ift_mask_fix_1N4n_9d30cad5_bbh_cot_fewshot_scatter_only.pdf}
    \includegraphics[width=0.554\textwidth, trim=0cm 0cm 0cm 0cm, clip]{auto_figures/dynamics_daint_prod_ift_q3-4b_1N4n_16cdce0f_bbh_cot_fewshot_scatter_only.pdf}
    \caption{The performance of our (\textbf{Left}) L3.1-8B-Magpie based MTP LM and (\textbf{Right}) Qwen3-4B-Inst-2507 MTP LM evaluated on the BBH benchmark after $\sim$100k steps of training. Performance tradeoff is visualized by plotting the effective $k$ value or ``Acceleration Factor'' versus the Accuracy on the benchmark. More detailed plots showing accuracy and acceleration as a function of training for both models are provided in \cref{fig:dynamics-flagship-l3-bbh,fig:dynamics-flagship-q3-bbh}. We observe that the adaptive decoding strategies achieve pareto-optimal tradeoffs between generation speed and response quality for both models. However, we note that for the Qwen3-4B-Inst-2507 model on this benchmark, the \textit{ConfAdapt} method appears to be similar to the \textit{Static} scheme at $k=3$ in terms of its performance speed tradeoff.}
    \label{fig:dynamics-flagship-l3-q3-bbh-scatter}
\end{figure*}
\begin{figure*}[h!]
    \includegraphics[width=\textwidth]{auto_figures/dynamics_daint_prod_ift_mask_fix_1N4n_9d30cad5_bbh_cot_fewshot_static_only.pdf}
    \includegraphics[width=\textwidth]{auto_figures/dynamics_daint_prod_ift_mask_fix_1N4n_9d30cad5_bbh_cot_fewshot_ca_only.pdf}
    \caption{(Left) Accuracy and (middle) acceleration dynamics for L3.1-8B-Magpie evaluated on BBH over the course of training. (Right) The pareto frontier of accuracy versus acceleration based on the final model checkpoint evaluated across all decoding strategies. We find that performance and acceleration stabilize around 50k training steps, and the adaptive decoding strategies achieve pareto-optimal tradeoffs between generation speed and response quality.}
    \label{fig:dynamics-flagship-l3-bbh}
\end{figure*}
\begin{figure*}[h!]
    \includegraphics[width=\textwidth]{auto_figures/dynamics_daint_prod_ift_q3-4b_1N4n_16cdce0f_bbh_cot_fewshot_static_only.pdf}
    \includegraphics[width=\textwidth]{auto_figures/dynamics_daint_prod_ift_q3-4b_1N4n_16cdce0f_bbh_cot_fewshot_ca_only.pdf}
    \caption{(Left) Accuracy and (middle) acceleration dynamics for Qwen3-4B-Inst-2507 evaluated on BBH over the course of training. (Right) The pareto frontier of accuracy versus acceleration based on the final model checkpoint evaluated across all decoding strategies. We find that performance and acceleration stabilize around 50k training steps, and the adaptive decoding strategies achieve pareto-optimal tradeoffs between generation speed and response quality.}
    \label{fig:dynamics-flagship-q3-bbh}
\end{figure*}

\clearpage

\subsection{Performance across tasks}
We also evaluate the same two models from \cref{fig:dynamics-flagship-l3-q3-gsm-scatter} on a more general set of evaluations, and summarize the performance of both models on both benchmarks in \cref{tab:math-evals} and \cref{tab:general-evals}.
Here we report the accuracy and acceleration metrics (``Effective $k$'' value) for the final checkpoint in the training run while also including the performance of each model at step 0 evaluated under standard NTP decoding (the pretrained initial checkpoint under the static $k=1$ scheme). 

We observe that even though both models were trained on only MetaMathQA, transfer learning occurs. For the L3.1-8B-Magpie model, in many cases the ConfAdapt scheme also achieves $>3\times$ acceleration with minor impact on performance relative to $k=1$ decoding. When compared to the performance of the original model checkpoints at step 0 (``Baseline''), the impact on accuracy appears more severe, but as these are highly optimized post-trained models that we finetune on a small dataset for a limited number of iterations, we think that this decrease even under $k=1$ decoding is unsurprising. 

As the two initial models have different strengths represented by performance at step 0 before MTP adaptation, we generally observe that larger average acceleration factors occur on the GSM8K and BBH benchmarks, and that for the more open-ended generative tasks like CNN DailyMail, relatively low acceleration factors are yielded for all adaptive decoding configurations.

% \begin{table*}[h!]
% \small
% \caption{Complete evaluation results for the two main models on the math and STEM centric benchmarks.}
\begin{table*}[h!]
\small
\caption{Complete evaluation results for the two main models on the math and STEM centric benchmarks.}
\label{tab:math-evals}
\centering
\begin{tabular}{llcccccc}
\toprule
 &  & \multicolumn{2}{c}{GSM8K} & \multicolumn{2}{c}{AIME25} & \multicolumn{2}{c}{GPQA Main} \\
 &  & Acc. (\%) & Eff. k & Acc. (\%) & Eff. k & Acc. (\%) & Eff. k \\
\midrule
L3.1-8B-Magpie & Baseline Step 0, k=1 & 69.5 ± 1.3 & 1 & 0.0 ± 0.0 & 1 & 17.6 ± 1.8 & 1 \\
 Chat Template: On & Static k=1 & 66.0 ± 1.3 & 1 & 0.0 ± 0.0 & 1 & 15.4 ± 1.7 & 1 \\
 & Static k=2 & 60.7 ± 1.3 & 2 & 0.0 ± 0.0 & 2 & 16.1 ± 1.7 & 2 \\
 & Static k=3 & 53.0 ± 1.4 & 3 & 0.0 ± 0.0 & 3 & 11.4 ± 1.5 & 3 \\
 & Static k=4 & 49.8 ± 1.4 & 4 & 0.0 ± 0.0 & 4 & 8.7 ± 1.3 & 4 \\
 & Static k=5 & 41.5 ± 1.4 & 5 & 0.0 ± 0.0 & 5 & 6.5 ± 1.2 & 5 \\
 & ConfAdapt ($\tau=0.995$) & 65.7 ± 1.3 & 1.9 ± 0.5 & 0.0 ± 0.0 & 3.0 ± 1.0 & 14.7 ± 1.7 & 1.5 ± 0.3 \\
 & ConfAdapt ($\tau=0.99$) & 66.0 ± 1.3 & 2.1 ± 0.6 & 0.0 ± 0.0 & 3.6 ± 1.4 & 14.3 ± 1.7 & 1.6 ± 0.3 \\
 & ConfAdapt ($\tau=0.98$) & 65.5 ± 1.3 & 2.3 ± 0.7 & 0.0 ± 0.0 & 4.3 ± 1.5 & 14.5 ± 1.7 & 1.7 ± 0.4 \\
 & ConfAdapt ($\tau=0.97$) & 65.3 ± 1.3 & 2.5 ± 0.8 & 0.0 ± 0.0 & 4.5 ± 1.8 & 14.3 ± 1.7 & 1.7 ± 0.4 \\
 & ConfAdapt ($\tau=0.96$) & 64.9 ± 1.3 & 2.6 ± 0.8 & 0.0 ± 0.0 & 4.7 ± 1.9 & 14.3 ± 1.7 & 1.8 ± 0.4 \\
 & ConfAdapt ($\tau=0.95$) & 64.7 ± 1.3 & 2.8 ± 0.9 & 0.0 ± 0.0 & 4.8 ± 2.0 & 15.2 ± 1.7 & 1.8 ± 0.4 \\
 & ConfAdapt ($\tau=0.9$) & 64.1 ± 1.3 & 3.3 ± 1.1 & 0.0 ± 0.0 & 5.5 ± 2.4 & 15.0 ± 1.7 & 2.0 ± 0.5 \\
 & ConfAdapt ($\tau=0.87$) & 62.3 ± 1.3 & 3.6 ± 1.2 & 0.0 ± 0.0 & 5.1 ± 2.2 & 14.7 ± 1.7 & 2.1 ± 0.5 \\
 & ConfAdapt ($\tau=0.85$) & 61.2 ± 1.3 & 3.7 ± 1.3 & 0.0 ± 0.0 & 5.2 ± 2.6 & 15.6 ± 1.7 & 2.2 ± 0.5 \\
 & ConfAdapt ($\tau=0.8$) & 59.2 ± 1.4 & 4.2 ± 1.5 & 0.0 ± 0.0 & 5.8 ± 2.6 & 15.6 ± 1.7 & 2.4 ± 0.6 \\
 & ConfAdapt ($\tau=0.75$) & 56.9 ± 1.4 & 4.6 ± 1.6 & 0.0 ± 0.0 & 5.7 ± 2.7 & 14.1 ± 1.6 & 2.5 ± 0.6 \\
 & ConfAdapt ($\tau=0.7$) & 55.0 ± 1.4 & 5.1 ± 1.8 & 0.0 ± 0.0 & 6.4 ± 3.2 & 17.2 ± 1.8 & 2.7 ± 0.6 \\
 & ConfAdapt ($\tau=0.65$) & 51.5 ± 1.4 & 5.6 ± 1.9 & 0.0 ± 0.0 & 5.7 ± 3.0 & 15.0 ± 1.7 & 2.9 ± 0.7 \\
 & ConfAdapt ($\tau=0.6$) & 46.6 ± 1.4 & 6.1 ± 2.1 & 0.0 ± 0.0 & 6.4 ± 3.3 & 14.1 ± 1.6 & 3.1 ± 0.7 \\
\cline{1-8}
Qwen3-4B-Inst-2507 & Baseline Step 0, k=1 & 75.4 ± 1.2 & 1 & 46.7 ± 9.3 & 1 & 12.3 ± 1.6 & 1 \\
 Chat Template: Off & Static k=1 & 89.1 ± 0.9 & 1 & 23.3 ± 7.9 & 1 & 15.8 ± 1.7 & 1 \\
 & Static k=2 & 82.3 ± 1.1 & 2 & 0.0 ± 0.0 & 2 & 18.3 ± 1.8 & 2 \\
 & Static k=3 & 61.6 ± 1.3 & 3 & 0.0 ± 0.0 & 3 & 7.1 ± 1.2 & 3 \\
 & Static k=4 & 44.2 ± 1.4 & 4 & 0.0 ± 0.0 & 4 & 6.0 ± 1.1 & 4 \\
 & Static k=5 & 28.9 ± 1.2 & 5 & 0.0 ± 0.0 & 5 & 2.7 ± 0.8 & 5 \\
 & ConfAdapt ($\tau=0.995$) & 88.2 ± 0.9 & 2.1 ± 0.4 & 26.7 ± 8.2 & 2.3 ± 0.4 & 15.0 ± 1.7 & 1.2 ± 0.1 \\
 & ConfAdapt ($\tau=0.99$) & 87.6 ± 0.9 & 2.2 ± 0.5 & 13.3 ± 6.3 & 2.5 ± 0.5 & 14.1 ± 1.6 & 1.2 ± 0.1 \\
 & ConfAdapt ($\tau=0.98$) & 87.5 ± 0.9 & 2.4 ± 0.6 & 20.0 ± 7.4 & 2.6 ± 0.5 & 13.4 ± 1.6 & 1.3 ± 0.1 \\
 & ConfAdapt ($\tau=0.97$) & 87.5 ± 0.9 & 2.5 ± 0.6 & 13.3 ± 6.3 & 2.6 ± 0.5 & 13.6 ± 1.6 & 1.3 ± 0.1 \\
 & ConfAdapt ($\tau=0.96$) & 87.0 ± 0.9 & 2.6 ± 0.7 & 16.7 ± 6.9 & 2.6 ± 0.6 & 13.2 ± 1.6 & 1.3 ± 0.1 \\
 & ConfAdapt ($\tau=0.95$) & 86.9 ± 0.9 & 2.7 ± 0.7 & 23.3 ± 7.9 & 2.8 ± 0.5 & 13.6 ± 1.6 & 1.3 ± 0.2 \\
 & ConfAdapt ($\tau=0.9$) & 83.6 ± 1.0 & 3.1 ± 0.8 & 10.0 ± 5.6 & 3.1 ± 0.7 & 14.7 ± 1.7 & 1.4 ± 0.2 \\
 & ConfAdapt ($\tau=0.87$) & 82.3 ± 1.1 & 3.2 ± 0.9 & 10.0 ± 5.6 & 3.2 ± 0.5 & 15.4 ± 1.7 & 1.5 ± 0.2 \\
 & ConfAdapt ($\tau=0.85$) & 81.3 ± 1.1 & 3.3 ± 0.9 & 6.7 ± 4.6 & 3.3 ± 0.6 & 14.3 ± 1.7 & 1.5 ± 0.2 \\
 & ConfAdapt ($\tau=0.8$) & 77.6 ± 1.1 & 3.6 ± 1.0 & 10.0 ± 5.6 & 3.4 ± 0.7 & 17.2 ± 1.8 & 1.6 ± 0.2 \\
 & ConfAdapt ($\tau=0.75$) & 73.2 ± 1.2 & 3.8 ± 1.1 & 6.7 ± 4.6 & 3.6 ± 0.9 & 15.8 ± 1.7 & 1.7 ± 0.3 \\
 & ConfAdapt ($\tau=0.7$) & 68.2 ± 1.3 & 4.1 ± 1.2 & 6.7 ± 4.6 & 3.4 ± 0.9 & 14.5 ± 1.7 & 1.8 ± 0.3 \\
 & ConfAdapt ($\tau=0.65$) & 62.0 ± 1.3 & 4.4 ± 1.3 & 3.3 ± 3.3 & 3.8 ± 1.4 & 15.8 ± 1.7 & 1.8 ± 0.3 \\
 & ConfAdapt ($\tau=0.6$) & 58.8 ± 1.4 & 4.6 ± 1.3 & 6.7 ± 4.6 & 4.2 ± 1.4 & 18.1 ± 1.8 & 2.0 ± 0.3 \\
\cline{1-8}
\bottomrule
\end{tabular}
\end{table*}




% \begin{table*}[h!]
% \small
% \caption{Complete evaluation results for the two main models on the instruction following and open ended generation benchmarks.}
\begin{table*}[h!]
\small
\caption{Complete evaluation results for the two main models on the instruction following and open ended generation benchmarks.}
\label{tab:general-evals}
\centering
\begin{tabular}{llcccccc}
\toprule
 &  & \multicolumn{2}{c}{BBH} & \multicolumn{2}{c}{IFEval} & \multicolumn{2}{c}{CNN DailyMail} \\
 &  & Acc. (\%) & Eff. k & Acc. (\%) & Eff. k & ROUGE-L & Eff. k \\
\midrule
L3.1-8B-Magpie & Baseline Step 0, k=1 & 67.1 ± 0.5 & 1 & 55.1 ± 2.1 & 1 & 0.226 ± 0.001 & 1 \\
 Chat Template: On & Static k=1 & 57.6 ± 0.5 & 1 & 42.5 ± 2.1 & 1 & 0.236 ± 0.001 & 1 \\
 & Static k=2 & 51.4 ± 0.5 & 2 & 29.8 ± 2.0 & 2 & 0.226 ± 0.001 & 2 \\
 & Static k=3 & 45.1 ± 0.5 & 3 & 23.5 ± 1.8 & 3 & 0.206 ± 0.001 & 3 \\
 & Static k=4 & 28.0 ± 0.5 & 4 & 22.0 ± 1.8 & 4 & 0.176 ± 0.001 & 4 \\
 & Static k=5 & 16.6 ± 0.4 & 5 & 19.8 ± 1.7 & 5 & 0.149 ± 0.001 & 5 \\
 & ConfAdapt ($\tau=0.995$) & 57.3 ± 0.5 & 2.4 ± 0.9 & 42.0 ± 2.1 & 1.5 ± 0.7 & 0.243 ± 0.001 & 1.5 ± 0.4 \\
 & ConfAdapt ($\tau=0.99$) & 56.8 ± 0.5 & 2.6 ± 1.0 & 41.8 ± 2.1 & 1.6 ± 0.8 & 0.243 ± 0.001 & 1.6 ± 0.5 \\
 & ConfAdapt ($\tau=0.98$) & 56.4 ± 0.5 & 3.0 ± 1.2 & 42.1 ± 2.1 & 1.8 ± 1.0 & 0.242 ± 0.001 & 1.7 ± 0.6 \\
 & ConfAdapt ($\tau=0.97$) & 56.2 ± 0.5 & 3.2 ± 1.3 & 42.5 ± 2.1 & 1.8 ± 1.1 & 0.242 ± 0.001 & 1.8 ± 0.6 \\
 & ConfAdapt ($\tau=0.96$) & 56.0 ± 0.5 & 3.3 ± 1.4 & 40.9 ± 2.1 & 1.9 ± 1.1 & 0.242 ± 0.001 & 1.8 ± 0.7 \\
 & ConfAdapt ($\tau=0.95$) & 55.6 ± 0.5 & 3.4 ± 1.4 & 41.4 ± 2.1 & 2.0 ± 1.2 & 0.242 ± 0.001 & 1.9 ± 0.7 \\
 & ConfAdapt ($\tau=0.9$) & 53.8 ± 0.5 & 3.9 ± 1.6 & 40.3 ± 2.1 & 2.2 ± 1.3 & 0.241 ± 0.001 & 2.1 ± 0.8 \\
 & ConfAdapt ($\tau=0.87$) & 51.9 ± 0.5 & 4.1 ± 1.7 & 40.5 ± 2.1 & 2.3 ± 1.4 & 0.241 ± 0.001 & 2.2 ± 0.8 \\
 & ConfAdapt ($\tau=0.85$) & 50.4 ± 0.5 & 4.2 ± 1.7 & 39.7 ± 2.1 & 2.3 ± 1.5 & 0.24 ± 0.001 & 2.2 ± 0.9 \\
 & ConfAdapt ($\tau=0.8$) & 45.5 ± 0.5 & 4.4 ± 1.8 & 38.6 ± 2.1 & 2.4 ± 1.5 & 0.239 ± 0.001 & 2.4 ± 0.9 \\
 & ConfAdapt ($\tau=0.75$) & 39.4 ± 0.5 & 4.6 ± 1.9 & 33.8 ± 2.0 & 2.6 ± 1.6 & 0.237 ± 0.001 & 2.5 ± 0.9 \\
 & ConfAdapt ($\tau=0.7$) & 33.6 ± 0.5 & 4.8 ± 2.0 & 35.5 ± 2.1 & 2.9 ± 1.9 & 0.236 ± 0.001 & 2.7 ± 1.0 \\
 & ConfAdapt ($\tau=0.65$) & 27.6 ± 0.5 & 5.1 ± 2.1 & 34.0 ± 2.0 & 3.1 ± 2.0 & 0.232 ± 0.001 & 2.9 ± 1.0 \\
 & ConfAdapt ($\tau=0.6$) & 21.2 ± 0.5 & 5.3 ± 2.2 & 32.2 ± 2.0 & 3.2 ± 2.0 & 0.229 ± 0.001 & 3.1 ± 1.0 \\
\cline{1-8}
Qwen3-4B-Inst-2507 & Baseline Step 0, k=1 & 34.4 ± 0.4 & 1 & 65.6 ± 2.0 & 1 & 0.223 ± 0.001 & 1 \\
 Chat Template: Off & Static k=1 & 28.0 ± 0.4 & 1 & 66.2 ± 2.0 & 1 & 0.227 ± 0.001 & 1 \\
 & Static k=2 & 26.9 ± 0.4 & 2 & 42.9 ± 2.1 & 2 & 0.222 ± 0.001 & 2 \\
 & Static k=3 & 21.0 ± 0.4 & 3 & 31.1 ± 2.0 & 3 & 0.205 ± 0.001 & 3 \\
 & Static k=4 & 14.9 ± 0.4 & 4 & 26.6 ± 1.9 & 4 & 0.177 ± 0.001 & 4 \\
 & Static k=5 & 10.2 ± 0.3 & 5 & 22.6 ± 1.8 & 5 & 0.156 ± 0.001 & 5 \\
 & ConfAdapt ($\tau=0.995$) & 27.4 ± 0.3 & 2.1 ± 0.8 & 64.7 ± 2.1 & 1.5 ± 0.5 & 0.235 ± 0.001 & 1.2 ± 0.1 \\
 & ConfAdapt ($\tau=0.99$) & 27.2 ± 0.3 & 2.3 ± 0.9 & 64.3 ± 2.1 & 1.5 ± 0.5 & 0.235 ± 0.001 & 1.2 ± 0.1 \\
 & ConfAdapt ($\tau=0.98$) & 26.8 ± 0.4 & 2.5 ± 1.0 & 62.3 ± 2.1 & 1.6 ± 0.6 & 0.234 ± 0.001 & 1.2 ± 0.1 \\
 & ConfAdapt ($\tau=0.97$) & 26.4 ± 0.4 & 2.6 ± 1.1 & 61.9 ± 2.1 & 1.7 ± 0.6 & 0.234 ± 0.001 & 1.3 ± 0.1 \\
 & ConfAdapt ($\tau=0.96$) & 26.2 ± 0.4 & 2.6 ± 1.1 & 62.1 ± 2.1 & 1.7 ± 0.7 & 0.234 ± 0.001 & 1.3 ± 0.1 \\
 & ConfAdapt ($\tau=0.95$) & 26.2 ± 0.4 & 2.7 ± 1.1 & 61.6 ± 2.1 & 1.7 ± 0.7 & 0.234 ± 0.001 & 1.3 ± 0.1 \\
 & ConfAdapt ($\tau=0.9$) & 26.1 ± 0.4 & 3.0 ± 1.2 & 59.5 ± 2.1 & 1.9 ± 0.9 & 0.233 ± 0.001 & 1.4 ± 0.1 \\
 & ConfAdapt ($\tau=0.87$) & 25.8 ± 0.4 & 3.1 ± 1.3 & 58.2 ± 2.1 & 2.0 ± 0.9 & 0.232 ± 0.001 & 1.4 ± 0.2 \\
 & ConfAdapt ($\tau=0.85$) & 25.2 ± 0.4 & 3.1 ± 1.3 & 59.3 ± 2.1 & 2.1 ± 1.0 & 0.232 ± 0.001 & 1.5 ± 0.2 \\
 & ConfAdapt ($\tau=0.8$) & 24.6 ± 0.4 & 3.3 ± 1.4 & 55.5 ± 2.1 & 2.2 ± 1.0 & 0.231 ± 0.001 & 1.5 ± 0.2 \\
 & ConfAdapt ($\tau=0.75$) & 23.5 ± 0.4 & 3.4 ± 1.4 & 55.3 ± 2.1 & 2.3 ± 1.1 & 0.23 ± 0.001 & 1.6 ± 0.2 \\
 & ConfAdapt ($\tau=0.7$) & 22.5 ± 0.4 & 3.6 ± 1.4 & 56.7 ± 2.1 & 2.5 ± 1.2 & 0.229 ± 0.001 & 1.7 ± 0.2 \\
 & ConfAdapt ($\tau=0.65$) & 21.8 ± 0.4 & 3.7 ± 1.5 & 51.8 ± 2.2 & 2.7 ± 1.3 & 0.228 ± 0.001 & 1.8 ± 0.3 \\
 & ConfAdapt ($\tau=0.6$) & 19.7 ± 0.4 & 4.0 ± 1.5 & 49.9 ± 2.2 & 3.0 ± 1.6 & 0.227 ± 0.001 & 1.9 ± 0.3 \\
\cline{1-8}
\bottomrule
\end{tabular}
\end{table*}



\clearpage

\subsection{Ablations}\label{app:ablations}

\paragraph{Alternate finetuning dataset.}
We also evaluate a version of the L3.1-8B-Magpie model further finetuned on the Magpie dataset instead of the MetaMathQA data used in all other experiments to analyze the impact of minimizing the distribution shift between the online MTP objective and the most recent stage of training for the initial checkpoint we use. In \cref{tab:abl-l3-data-evals} we can see that while, the $k=1$ performance of the models is quite similar, and there are a few instances of decoding strategy and task pairings where the model trained on Magpie performs better, in general, the version trained on MetaMathQA achieves better accuracy and acceleration in most cases. We suspect that this is due to the fact that the Magpie datasets are much larger and therefore at the same iteration count

% \begin{table*}[h!]
% \small
% \caption{Complete evaluation results for the main L3.1-8B-Magpie model trained on MetaMathQA versus the version trained on Magpie on a subset of the evaluation benchmarks.}
\begin{table*}[h!]
\small
\caption{Complete evaluation results for the main L3.1-8B-Magpie model trained on MetaMathQA versus the version trained on Magpie on a subset of the evaluation benchmarks.}
\label{tab:abl-l3-data-evals}
\centering
\begin{tabular}{llcccccc}
\toprule
 &  & \multicolumn{2}{c}{GSM8K} & \multicolumn{2}{c}{BBH} & \multicolumn{2}{c}{IFEval} \\
 &  & Acc. (\%) & Eff. k & Acc. (\%) & Eff. k & Acc. (\%) & Eff. k \\
\midrule
L3.1-8B-Magpie & Baseline Step 0, k=1 & 69.5 ± 1.3 & 1 & 67.1 ± 0.5 & 1 & 55.1 ± 2.1 & 1 \\
 Chat Template: On & Static k=1 & 66.0 ± 1.3 & 1 & 57.6 ± 0.5 & 1 & 42.5 ± 2.1 & 1 \\
 & Static k=2 & 60.7 ± 1.3 & 2 & 51.4 ± 0.5 & 2 & 29.8 ± 2.0 & 2 \\
 & Static k=3 & 53.0 ± 1.4 & 3 & 45.1 ± 0.5 & 3 & 23.5 ± 1.8 & 3 \\
 & Static k=4 & 49.8 ± 1.4 & 4 & 28.0 ± 0.5 & 4 & 22.0 ± 1.8 & 4 \\
 & Static k=5 & 41.5 ± 1.4 & 5 & 16.6 ± 0.4 & 5 & 19.8 ± 1.7 & 5 \\
 & ConfAdapt ($\tau=0.995$) & 65.7 ± 1.3 & 1.9 ± 0.5 & 57.3 ± 0.5 & 2.4 ± 0.9 & 42.0 ± 2.1 & 1.5 ± 0.7 \\
 & ConfAdapt ($\tau=0.99$) & 66.0 ± 1.3 & 2.1 ± 0.6 & 56.8 ± 0.5 & 2.6 ± 1.0 & 41.8 ± 2.1 & 1.6 ± 0.8 \\
 & ConfAdapt ($\tau=0.98$) & 65.5 ± 1.3 & 2.3 ± 0.7 & 56.4 ± 0.5 & 3.0 ± 1.2 & 42.1 ± 2.1 & 1.8 ± 1.0 \\
 & ConfAdapt ($\tau=0.97$) & 65.3 ± 1.3 & 2.5 ± 0.8 & 56.2 ± 0.5 & 3.2 ± 1.3 & 42.5 ± 2.1 & 1.8 ± 1.1 \\
 & ConfAdapt ($\tau=0.96$) & 64.9 ± 1.3 & 2.6 ± 0.8 & 56.0 ± 0.5 & 3.3 ± 1.4 & 40.9 ± 2.1 & 1.9 ± 1.1 \\
 & ConfAdapt ($\tau=0.95$) & 64.7 ± 1.3 & 2.8 ± 0.9 & 55.6 ± 0.5 & 3.4 ± 1.4 & 41.4 ± 2.1 & 2.0 ± 1.2 \\
 & ConfAdapt ($\tau=0.9$) & 64.1 ± 1.3 & 3.3 ± 1.1 & 53.8 ± 0.5 & 3.9 ± 1.6 & 40.3 ± 2.1 & 2.2 ± 1.3 \\
 & ConfAdapt ($\tau=0.87$) & 62.3 ± 1.3 & 3.6 ± 1.2 & 51.9 ± 0.5 & 4.1 ± 1.7 & 40.5 ± 2.1 & 2.3 ± 1.4 \\
 & ConfAdapt ($\tau=0.85$) & 61.2 ± 1.3 & 3.7 ± 1.3 & 50.4 ± 0.5 & 4.2 ± 1.7 & 39.7 ± 2.1 & 2.3 ± 1.5 \\
 & ConfAdapt ($\tau=0.8$) & 59.2 ± 1.4 & 4.2 ± 1.5 & 45.5 ± 0.5 & 4.4 ± 1.8 & 38.6 ± 2.1 & 2.4 ± 1.5 \\
 & ConfAdapt ($\tau=0.75$) & 56.9 ± 1.4 & 4.6 ± 1.6 & 39.4 ± 0.5 & 4.6 ± 1.9 & 33.8 ± 2.0 & 2.6 ± 1.6 \\
 & ConfAdapt ($\tau=0.7$) & 55.0 ± 1.4 & 5.1 ± 1.8 & 33.6 ± 0.5 & 4.8 ± 2.0 & 35.5 ± 2.1 & 2.9 ± 1.9 \\
 & ConfAdapt ($\tau=0.65$) & 51.5 ± 1.4 & 5.6 ± 1.9 & 27.6 ± 0.5 & 5.1 ± 2.1 & 34.0 ± 2.0 & 3.1 ± 2.0 \\
 & ConfAdapt ($\tau=0.6$) & 46.6 ± 1.4 & 6.1 ± 2.1 & 21.2 ± 0.5 & 5.3 ± 2.2 & 32.2 ± 2.0 & 3.2 ± 2.0 \\
\cline{1-8}
L3.1-8B-Magpie FT Magpie & Baseline Step 0, k=1 & 69.5 ± 1.3 & 1 & 67.1 ± 0.5 & 1 & 55.1 ± 2.1 & 1 \\
 Chat Template: On & Static k=1 & 64.5 ± 1.3 & 1 & 57.7 ± 0.5 & 1 & 41.8 ± 2.1 & 1 \\
 & Static k=2 & 53.8 ± 1.4 & 2 & 54.2 ± 0.5 & 2 & 30.3 ± 2.0 & 2 \\
 & Static k=3 & 45.3 ± 1.4 & 3 & 49.9 ± 0.5 & 3 & 27.9 ± 1.9 & 3 \\
 & Static k=4 & 34.0 ± 1.3 & 4 & 34.2 ± 0.5 & 4 & 24.6 ± 1.9 & 4 \\
 & Static k=5 & 26.2 ± 1.2 & 5 & 24.9 ± 0.5 & 5 & 20.0 ± 1.7 & 5 \\
 & ConfAdapt ($\tau=0.995$) & 64.5 ± 1.3 & 1.5 ± 0.3 & 57.3 ± 0.5 & 2.8 ± 1.1 & 41.6 ± 2.1 & 1.5 ± 0.9 \\
 & ConfAdapt ($\tau=0.99$) & 64.6 ± 1.3 & 1.7 ± 0.3 & 54.9 ± 0.5 & 3.2 ± 1.4 & 41.8 ± 2.1 & 1.6 ± 1.1 \\
 & ConfAdapt ($\tau=0.98$) & 64.4 ± 1.3 & 1.8 ± 0.4 & 46.7 ± 0.5 & 3.7 ± 1.7 & 41.0 ± 2.1 & 1.7 ± 1.3 \\
 & ConfAdapt ($\tau=0.97$) & 64.0 ± 1.3 & 2.0 ± 0.5 & 38.6 ± 0.5 & 4.0 ± 1.9 & 41.0 ± 2.1 & 1.7 ± 1.4 \\
 & ConfAdapt ($\tau=0.96$) & 64.1 ± 1.3 & 2.1 ± 0.5 & 34.6 ± 0.5 & 4.3 ± 2.0 & 40.7 ± 2.1 & 1.8 ± 1.5 \\
 & ConfAdapt ($\tau=0.95$) & 63.5 ± 1.3 & 2.1 ± 0.6 & 31.0 ± 0.5 & 4.5 ± 2.1 & 40.5 ± 2.1 & 1.8 ± 1.5 \\
 & ConfAdapt ($\tau=0.9$) & 61.5 ± 1.3 & 2.5 ± 0.7 & 18.8 ± 0.4 & 5.1 ± 2.4 & 39.7 ± 2.1 & 2.0 ± 1.6 \\
 & ConfAdapt ($\tau=0.87$) & 60.7 ± 1.3 & 2.7 ± 0.8 & 16.8 ± 0.4 & 5.4 ± 2.5 & 39.6 ± 2.1 & 2.1 ± 1.8 \\
 & ConfAdapt ($\tau=0.85$) & 59.4 ± 1.4 & 2.8 ± 0.8 & 15.8 ± 0.4 & 5.6 ± 2.5 & 39.4 ± 2.1 & 2.1 ± 1.8 \\
 & ConfAdapt ($\tau=0.8$) & 54.9 ± 1.4 & 3.1 ± 0.9 & 14.7 ± 0.4 & 6.0 ± 2.7 & 38.8 ± 2.1 & 2.3 ± 1.9 \\
 & ConfAdapt ($\tau=0.75$) & 50.9 ± 1.4 & 3.4 ± 1.0 & 13.2 ± 0.4 & 6.3 ± 2.7 & 37.7 ± 2.1 & 2.4 ± 2.0 \\
 & ConfAdapt ($\tau=0.7$) & 47.9 ± 1.4 & 3.7 ± 1.1 & 12.3 ± 0.4 & 6.7 ± 2.8 & 36.4 ± 2.1 & 2.6 ± 2.1 \\
 & ConfAdapt ($\tau=0.65$) & 42.6 ± 1.4 & 4.0 ± 1.2 & 10.2 ± 0.3 & 7.1 ± 2.9 & 35.7 ± 2.1 & 2.9 ± 2.3 \\
 & ConfAdapt ($\tau=0.6$) & 37.4 ± 1.3 & 4.3 ± 1.3 & 8.8 ± 0.3 & 7.4 ± 3.0 & 33.5 ± 2.0 & 3.2 ± 2.5 \\
\cline{1-8}
\bottomrule
\end{tabular}
\end{table*}


\clearpage

\paragraph{Bidirectional attention in MTP region.}

Due to the fact that the student MTP LM does not receive ground truth tokens as input within the masked region where it is performing prediction during training, our training objective does not have the potential for the ``anti-causal'' leakage that would occur in traditional offline NTP training if a \textit{bidirectional} attention mask was used. Thus, we are able to consider allowing all neighboring mask token representations within the same MTP region to attend to each other. As shown in \cref{tab:abl-l3-suprv-evals-pt1}, we find that increasing the representational and computational capacity of the MTP prediction pass does not yield a meaningfully better final model under the same training conditions.

\paragraph{Prefix NTP supervision}

Since we prepare input data during training using the $M=N/(2k)$ strided convention, as opposed to say $M=N/k$, certain regions in the sequence are ``prefix'' only (see ``Pred'' regions marked false/0 in \cref{fig:explainer-tok-masking}). Due to the carefully blocked attention mask, the logits at those positions in the MTP model's output are produced by a computation that exactly matches the standard NTP forward pass. As a result, we consider also supervising the standard NTP objective on all of the prefix positions but as shown in \cref{tab:abl-l3-suprv-evals-pt2}, we find that it underperforms the objective configuration used for the main experiments.

\clearpage

% \begin{table*}[h!]
% \small
% \caption{Evaluation results for the ablations of the training objective configuration on the L3.1-8B-Magpie model measured on GSM8K. The first group labeled just \textbf{``L3.1-8B-Magpie''} is the configuration used for the main experiments, the \textbf{``GT Suprv.''} uses the ground truth tokens as labels rather than the online teacher forcing, \textbf{``BDA''} allows bidirectional attention in the MTP regions. table continues below.}
\begin{table*}[h!]
\small
\caption{Evaluation results for the ablations of the training objective configuration on the L3.1-8B-Magpie model measured on GSM8K. The first group labeled just \textbf{``L3.1-8B-Magpie''} is the configuration used for the main experiments, the \textbf{``GT Suprv.''} uses the ground truth tokens as labels rather than the online teacher forcing, \textbf{``BDA''} allows bidirectional attention in the MTP regions. table continues below.}
\label{tab:abl-l3-suprv-evals-pt1}
\centering
\begin{tabular}{llcc}
\toprule
 &  & \multicolumn{2}{c}{GSM8K} \\
 &  & Acc. (\%) & Eff. k \\
\midrule
L3.1-8B-Magpie & Baseline Step 0, k=1 & 69.5 ± 1.3 & 1 \\
 Chat Template: On & Static k=1 & 66.0 ± 1.3 & 1 \\
 & Static k=2 & 60.7 ± 1.3 & 2 \\
 & Static k=3 & 53.0 ± 1.4 & 3 \\
 & Static k=4 & 49.8 ± 1.4 & 4 \\
 & Static k=5 & 41.5 ± 1.4 & 5 \\
 & ConfAdapt ($\tau=0.995$) & 65.7 ± 1.3 & 1.9 ± 0.5 \\
 & ConfAdapt ($\tau=0.99$) & 66.0 ± 1.3 & 2.1 ± 0.6 \\
 & ConfAdapt ($\tau=0.98$) & 65.5 ± 1.3 & 2.3 ± 0.7 \\
 & ConfAdapt ($\tau=0.97$) & 65.3 ± 1.3 & 2.5 ± 0.8 \\
 & ConfAdapt ($\tau=0.96$) & 64.9 ± 1.3 & 2.6 ± 0.8 \\
 & ConfAdapt ($\tau=0.95$) & 64.7 ± 1.3 & 2.8 ± 0.9 \\
 & ConfAdapt ($\tau=0.9$) & 64.1 ± 1.3 & 3.3 ± 1.1 \\
 & ConfAdapt ($\tau=0.87$) & 62.3 ± 1.3 & 3.6 ± 1.2 \\
 & ConfAdapt ($\tau=0.85$) & 61.2 ± 1.3 & 3.7 ± 1.3 \\
 & ConfAdapt ($\tau=0.8$) & 59.2 ± 1.4 & 4.2 ± 1.5 \\
 & ConfAdapt ($\tau=0.75$) & 56.9 ± 1.4 & 4.6 ± 1.6 \\
 & ConfAdapt ($\tau=0.7$) & 55.0 ± 1.4 & 5.1 ± 1.8 \\
 & ConfAdapt ($\tau=0.65$) & 51.5 ± 1.4 & 5.6 ± 1.9 \\
 & ConfAdapt ($\tau=0.6$) & 46.6 ± 1.4 & 6.1 ± 2.1 \\
\cline{1-4}
L3.1-8B-Magpie GT Suprv. & Baseline Step 0, k=1 & 69.5 ± 1.3 & 1 \\
 Chat Template: On & Static k=1 & 67.9 ± 1.3 & 1 \\
 & Static k=2 & 59.3 ± 1.4 & 2 \\
 & Static k=3 & 42.9 ± 1.4 & 3 \\
 & Static k=4 & 27.8 ± 1.2 & 4 \\
 & Static k=5 & 19.0 ± 1.1 & 5 \\
 & ConfAdapt ($\tau=0.995$) & 67.7 ± 1.3 & 1.4 ± 0.2 \\
 & ConfAdapt ($\tau=0.99$) & 67.6 ± 1.3 & 1.5 ± 0.3 \\
 & ConfAdapt ($\tau=0.98$) & 67.5 ± 1.3 & 1.7 ± 0.3 \\
 & ConfAdapt ($\tau=0.97$) & 67.6 ± 1.3 & 1.8 ± 0.4 \\
 & ConfAdapt ($\tau=0.96$) & 67.5 ± 1.3 & 1.9 ± 0.4 \\
 & ConfAdapt ($\tau=0.95$) & 67.3 ± 1.3 & 1.9 ± 0.4 \\
 & ConfAdapt ($\tau=0.9$) & 66.4 ± 1.3 & 2.2 ± 0.6 \\
 & ConfAdapt ($\tau=0.87$) & 66.2 ± 1.3 & 2.4 ± 0.6 \\
 & ConfAdapt ($\tau=0.85$) & 65.9 ± 1.3 & 2.5 ± 0.7 \\
 & ConfAdapt ($\tau=0.8$) & 64.3 ± 1.3 & 2.7 ± 0.8 \\
 & ConfAdapt ($\tau=0.75$) & 62.0 ± 1.3 & 2.9 ± 0.8 \\
 & ConfAdapt ($\tau=0.7$) & 59.0 ± 1.4 & 3.2 ± 0.9 \\
 & ConfAdapt ($\tau=0.65$) & 54.2 ± 1.4 & 3.4 ± 1.0 \\
 & ConfAdapt ($\tau=0.6$) & 49.9 ± 1.4 & 3.7 ± 1.1 \\
\cline{1-4}
L3.1-8B-Magpie BDA & Baseline Step 0, k=1 & 69.5 ± 1.3 & 1 \\
 Chat Template: On & Static k=1 & 63.5 ± 1.3 & 1 \\
 & Static k=2 & 56.0 ± 1.4 & 2 \\
 & Static k=3 & 47.9 ± 1.4 & 3 \\
 & Static k=4 & 43.1 ± 1.4 & 4 \\
 & Static k=5 & 37.2 ± 1.3 & 5 \\
 & ConfAdapt ($\tau=0.995$) & 63.4 ± 1.3 & 1.9 ± 0.5 \\
 & ConfAdapt ($\tau=0.99$) & 63.5 ± 1.3 & 2.1 ± 0.6 \\
 & ConfAdapt ($\tau=0.98$) & 63.3 ± 1.3 & 2.3 ± 0.7 \\
 & ConfAdapt ($\tau=0.97$) & 63.2 ± 1.3 & 2.4 ± 0.7 \\
 & ConfAdapt ($\tau=0.96$) & 62.8 ± 1.3 & 2.6 ± 0.8 \\
 & ConfAdapt ($\tau=0.95$) & 62.7 ± 1.3 & 2.7 ± 0.8 \\
 & ConfAdapt ($\tau=0.9$) & 59.8 ± 1.4 & 3.1 ± 1.0 \\
 & ConfAdapt ($\tau=0.87$) & 58.9 ± 1.4 & 3.3 ± 1.1 \\
 & ConfAdapt ($\tau=0.85$) & 58.5 ± 1.4 & 3.5 ± 1.2 \\
 & ConfAdapt ($\tau=0.8$) & 55.1 ± 1.4 & 3.8 ± 1.3 \\
 & ConfAdapt ($\tau=0.75$) & 52.2 ± 1.4 & 4.2 ± 1.4 \\
 & ConfAdapt ($\tau=0.7$) & 49.0 ± 1.4 & 4.6 ± 1.5 \\
 & ConfAdapt ($\tau=0.65$) & 46.0 ± 1.4 & 5.0 ± 1.6 \\
 & ConfAdapt ($\tau=0.6$) & 41.3 ± 1.4 & 5.3 ± 1.6 \\
\cline{1-4}
\bottomrule
\end{tabular}
\end{table*}



% \begin{table*}[h!]
% \small
% \caption{Evaluation results for the ablations of the training objective configuration on the L3.1-8B-Magpie model measured on GSM8K. Table continues from above. The \textbf{``Soft Teacher''} is trained using the full teacher logits rather than just the argmax as labels, \textbf{``Static $k$''} keeps $k$ fixed at the max value of 16 throughout training, and \textbf{``Prefix Loss''} includes the auxiliary loss term supervising the standard NTP loss on the non-MTP positions in each sequence.}
\begin{table*}[h!]
\small
\caption{Evaluation results for the ablations of the training objective configuration on the L3.1-8B-Magpie model measured on GSM8K. Table continues from above. The \textbf{``Soft Teacher''} is trained using the full teacher logits rather than just the argmax as labels, \textbf{``Static $k$''} keeps $k$ fixed at the max value of 16 throughout training, and \textbf{``Prefix Loss''} includes the auxiliary loss term supervising the standard NTP loss on the non-MTP positions in each sequence.}
\label{tab:abl-l3-suprv-evals-pt2}
\centering
\begin{tabular}{llcc}
\toprule
 &  & \multicolumn{2}{c}{GSM8K} \\
 &  & Acc. (\%) & Eff. k \\
\midrule
L3.1-8B-Magpie Soft Teacher & Baseline Step 0, k=1 & 69.5 ± 1.3 & 1 \\
 Chat Template: On & Static k=1 & 63.6 ± 1.3 & 1 \\
 & Static k=2 & 56.1 ± 1.4 & 2 \\
 & Static k=3 & 49.7 ± 1.4 & 3 \\
 & Static k=4 & 45.6 ± 1.4 & 4 \\
 & Static k=5 & 39.2 ± 1.3 & 5 \\
 & ConfAdapt ($\tau=0.995$) & 63.5 ± 1.3 & 1.3 ± 0.2 \\
 & ConfAdapt ($\tau=0.99$) & 63.5 ± 1.3 & 1.5 ± 0.3 \\
 & ConfAdapt ($\tau=0.98$) & 63.6 ± 1.3 & 1.7 ± 0.4 \\
 & ConfAdapt ($\tau=0.97$) & 63.6 ± 1.3 & 1.8 ± 0.5 \\
 & ConfAdapt ($\tau=0.96$) & 63.6 ± 1.3 & 1.9 ± 0.5 \\
 & ConfAdapt ($\tau=0.95$) & 63.5 ± 1.3 & 2.0 ± 0.5 \\
 & ConfAdapt ($\tau=0.9$) & 62.5 ± 1.3 & 2.3 ± 0.7 \\
 & ConfAdapt ($\tau=0.87$) & 62.1 ± 1.3 & 2.5 ± 0.8 \\
 & ConfAdapt ($\tau=0.85$) & 62.1 ± 1.3 & 2.6 ± 0.9 \\
 & ConfAdapt ($\tau=0.8$) & 61.0 ± 1.3 & 2.9 ± 1.0 \\
 & ConfAdapt ($\tau=0.75$) & 58.9 ± 1.4 & 3.2 ± 1.2 \\
 & ConfAdapt ($\tau=0.7$) & 56.8 ± 1.4 & 3.6 ± 1.3 \\
 & ConfAdapt ($\tau=0.65$) & 53.7 ± 1.4 & 4.0 ± 1.5 \\
 & ConfAdapt ($\tau=0.6$) & 49.3 ± 1.4 & 4.5 ± 1.8 \\
\cline{1-4}
L3.1-8B-Magpie Static k & Baseline Step 0, k=1 & 69.5 ± 1.3 & 1 \\
 Chat Template: On & Static k=1 & 54.1 ± 1.4 & 1 \\
 & Static k=2 & 45.6 ± 1.4 & 2 \\
 & Static k=3 & 38.7 ± 1.3 & 3 \\
 & Static k=4 & 35.1 ± 1.3 & 4 \\
 & Static k=5 & 32.4 ± 1.3 & 5 \\
 & ConfAdapt ($\tau=0.995$) & 54.1 ± 1.4 & 1.5 ± 0.3 \\
 & ConfAdapt ($\tau=0.99$) & 54.1 ± 1.4 & 1.7 ± 0.4 \\
 & ConfAdapt ($\tau=0.98$) & 54.1 ± 1.4 & 1.9 ± 0.6 \\
 & ConfAdapt ($\tau=0.97$) & 53.4 ± 1.4 & 2.0 ± 0.7 \\
 & ConfAdapt ($\tau=0.96$) & 53.2 ± 1.4 & 2.2 ± 0.8 \\
 & ConfAdapt ($\tau=0.95$) & 53.4 ± 1.4 & 2.3 ± 0.8 \\
 & ConfAdapt ($\tau=0.9$) & 51.6 ± 1.4 & 2.7 ± 1.1 \\
 & ConfAdapt ($\tau=0.87$) & 50.6 ± 1.4 & 3.0 ± 1.2 \\
 & ConfAdapt ($\tau=0.85$) & 50.6 ± 1.4 & 3.1 ± 1.3 \\
 & ConfAdapt ($\tau=0.8$) & 48.7 ± 1.4 & 3.5 ± 1.5 \\
 & ConfAdapt ($\tau=0.75$) & 45.9 ± 1.4 & 3.9 ± 1.6 \\
 & ConfAdapt ($\tau=0.7$) & 44.0 ± 1.4 & 4.3 ± 1.8 \\
 & ConfAdapt ($\tau=0.65$) & 41.5 ± 1.4 & 4.8 ± 2.0 \\
 & ConfAdapt ($\tau=0.6$) & 39.0 ± 1.3 & 5.3 ± 2.2 \\
\cline{1-4}
L3.1-8B-Magpie Prefix Loss & Baseline Step 0, k=1 & 69.5 ± 1.3 & 1 \\
 Chat Template: On & Static k=1 & 66.6 ± 1.3 & 1 \\
 & Static k=2 & 57.8 ± 1.4 & 2 \\
 & Static k=3 & 47.3 ± 1.4 & 3 \\
 & Static k=4 & 45.8 ± 1.4 & 4 \\
 & Static k=5 & 37.5 ± 1.3 & 5 \\
 & ConfAdapt ($\tau=0.995$) & 69.4 ± 1.3 & 1.6 ± 0.3 \\
 & ConfAdapt ($\tau=0.99$) & 69.5 ± 1.3 & 1.8 ± 0.4 \\
 & ConfAdapt ($\tau=0.98$) & 69.6 ± 1.3 & 2.0 ± 0.6 \\
 & ConfAdapt ($\tau=0.97$) & 69.3 ± 1.3 & 2.2 ± 0.7 \\
 & ConfAdapt ($\tau=0.96$) & 69.4 ± 1.3 & 2.3 ± 0.8 \\
 & ConfAdapt ($\tau=0.95$) & 69.3 ± 1.3 & 2.4 ± 0.8 \\
 & ConfAdapt ($\tau=0.9$) & 64.6 ± 1.3 & 2.7 ± 0.9 \\
 & ConfAdapt ($\tau=0.87$) & 66.3 ± 1.3 & 3.1 ± 1.1 \\
 & ConfAdapt ($\tau=0.85$) & 65.7 ± 1.3 & 3.3 ± 1.2 \\
 & ConfAdapt ($\tau=0.8$) & 63.2 ± 1.3 & 3.6 ± 1.3 \\
 & ConfAdapt ($\tau=0.75$) & 58.3 ± 1.4 & 4.0 ± 1.5 \\
 & ConfAdapt ($\tau=0.7$) & 54.4 ± 1.4 & 4.4 ± 1.6 \\
 & ConfAdapt ($\tau=0.65$) & 50.9 ± 1.4 & 4.9 ± 1.8 \\
 & ConfAdapt ($\tau=0.6$) & 47.5 ± 1.4 & 5.0 ± 1.7 \\
\cline{1-4}
\bottomrule
\end{tabular}
\end{table*}



\end{document}

% This document was modified from the file originally made available by
% Pat Langley and Andrea Danyluk for ICML-2K. This version was created
% by Iain Murray in 2018, and modified by Alexandre Bouchard in
% 2019 and 2021 and by Csaba Szepesvari, Gang Niu and Sivan Sabato in 2022.
% Modified again in 2023 and 2024 by Sivan Sabato and Jonathan Scarlett.
% Previous contributors include Dan Roy, Lise Getoor and Tobias
% Scheffer, which was slightly modified from the 2010 version by
% Thorsten Joachims & Johannes Fuernkranz, slightly modified from the
% 2009 version by Kiri Wagstaff and Sam Roweis's 2008 version, which is
% slightly modified from Prasad Tadepalli's 2007 version which is a
% lightly changed version of the previous year's version by Andrew
% Moore, which was in turn edited from those of Kristian Kersting and
% Codrina Lauth. Alex Smola contributed to the algorithmic style files.
